\section{Graphs with least eigenvalue \texorpdfstring{$-2$}{-2}}\label{sec:4}
\begingroup
\emergencystretch=3em

We will follow Section 4 of [CGSS]. The paper's graphs $G$ and $H$ are our graphs $X$ and $Y$.
All graphs are finite and simple. For a graph $X$ we write $A=A(X)$ for its adjacency matrix,
$J$ for the all-ones matrix and $\lambda_{\min}(X)$ for its least eigenvalue; the least
eigenvalue of $X$ is at least $-2$ exactly when $2I+A$ is positive semidefinite, and we will
state the results of this section in that form, as the Lean statements do. The coordinates
$\Ze_1,\dots,\Ze_\ell$ are those of Definition~\ref{def:3.1}. For a vector
$\Zx=\sum_i\Zx(i)\Ze_i$ the \emph{support} of $\Zx$ is $\supp(\Zx)=\{i:\Zx(i)\ne0\}$. For a graph $Y$,
the line graph $L(Y)$ has the edges of $Y$ as its vertices, two of them adjacent when they are
distinct and share an end.

\begingroup\raggedright
\medskip\noindent\textbf{Facts from Mathlib and the project.}\label{facts:sec4}
We will use the following declarations of Mathlib (commit \path{082e2d37}):
\path{Matrix.posSemidef_iff_eq_sum_vecMulVec} (a positive semidefinite real matrix is a sum of
matrices $\Zv\Zv^T$), \path{Matrix.posSemidef_gram} (a Gram matrix is positive semidefinite),
\path{Matrix.rank_eq_finrank_span_cols}, \path{Submodule.map_span},
\path{LinearMap.finrank_range_of_inj}, \path{Submodule.finrank_le}, \path{stdOrthonormalBasis}
with its field \path{OrthonormalBasis.repr},
\path{Matrix.det_gram_ne_zero_iff_linearIndependent},
\path{Matrix.posDef_gram_iff_linearIndependent},
\path{LinearIndependent.fintype_card_le_finrank}, \path{finrank_euclideanSpace_fin}, the
definition \path{Matrix.rank} (the dimension of the range), \path{Submodule.span_mono},
\path{Submodule.finrank_mono}, \path{Submodule.finrank_sup_add_finrank_inf_eq},
\path{finrank_span_singleton}, \path{SimpleGraph.adjMatrix_mulVec_const_apply_of_regular}, for a
real symmetric matrix its eigenvalues \path{Matrix.IsHermitian.eigenvalues} with
\path{Matrix.IsHermitian.spectral_theorem}, \path{Matrix.IsHermitian.trace_eq_sum_eigenvalues}
and \path{Matrix.IsHermitian.rank_eq_card_non_zero_eigs}, \path{Matrix.trace_mul_comm}, the
Cauchy--Schwarz inequality \path{Finset.sum_mul_sq_le_sq_mul_sq}, the predicate
\path{SimpleGraph.IsBipartite} (a 2-colouring) with
\path{SimpleGraph.two_colorable_iff_forall_loop_even} and
\path{SimpleGraph.Coloring.even_length_iff_congr}, and the structure
\path{SimpleGraph.IsSRGWith}. From the
project we use the bridge \path{SimpleGraph.posSemidef_two_add_adjMatrix_iff} ($2I+A$ is
positive semidefinite exactly when every eigenvalue of $X$ is at least $-2$) and
\path{SimpleGraph.leastEigenvalue_eq_neg_two_iff} (for a graph with a vertex,
$\lambda_{\min}(X)=-2$ exactly when $2I+A$ is positive semidefinite and singular), both in
\path{LineGraphs/Spectral/Positivity.lean}; and the sets
$N_\Zr=\{\Zx\in\Phi:(\Zx,\Zr)=1\}$ of \path{LineGraphs/Computations/NeighborCertificates.lean}.
\par
\endgroup

\subsection{Theorem 4.1}

\begin{namedresult}{Definition}{(representations)}\label{def:represented}
 \lean{SimpleGraph.IsRepresentation}\lean{SimpleGraph.IsRepresentedIn}
 \uses{def:line-system,def:3.2}
 \assume{$X$ is a graph, $V$ is a real inner product space and $S$ is a line system in $V$.}
 \pflet{A \emph{representation} of $X$ in $V$ is a family $(\Zx_u)_{u\in V(X)}$ of vectors of
 $V$ whose Gram matrix is $2I+A(X)$; that is, $(\Zx_u,\Zx_u)=2$, $(\Zx_u,\Zx_w)=1$ when $u$ and
 $w$ are adjacent, and $(\Zx_u,\Zx_w)=0$ when $u\ne w$ are not adjacent. The graph $X$ is
 \emph{represented in $S$} if it has a representation with $\Zx_u\in\Phi(S)$ for every $u$.}
 \new\ The paper speaks of a set of vectors and of the graph it represents; we will index the
 vectors by the vertices, so that the graph is given first, as in Lean. For $S=\typ E_8$ we take
 model (ii) of Definition~\ref{def:3.2} in $\re^8$.
\end{namedresult}

\begin{namedresult}{Lemma}{(Gram factorization)}\label{lem:4-gramfactor}
 \lean{Matrix.rank_gram_eq_finrank_span}\lean{Matrix.PosSemidef.exists_gram_euclideanSpace}
 \uses{facts:sec4}
 \assume{$M$ is a real symmetric matrix indexed by a finite set $V$, and $r\ge0$.}
 \prove{\begin{pfenum}
  \item If $M$ is the Gram matrix of a family $(\Zx_u)_{u\in V}$ in a real inner product space,
   then $\rk M$ is the dimension of the span of the $\Zx_u$.
  \item If $M$ is positive semidefinite and $\rk M\le r$, then $M$ is the Gram matrix of a
   family in $\re^r$, that is, in \texttt{EuclideanSpace}~$\re$~\texttt{(Fin r)}.
 \end{pfenum}}
 \new\ Mathlib factors a positive semidefinite matrix through a sum of an unspecified number of
 terms; the sections below need the dimension.
\end{namedresult}
\begin{proof}
 \pfsketch\ We will compare the columns of $M$ with the span, and then move the span into
 $\re^r$ by an isometry.
 \step{s4gfspan}{Assertion 1 holds.}
 \begin{proof}\pf\ Let $W$ be the span of the $\Zx_u$ and $T\Zy=((\Zx_w,\Zy))_{w\in V}$ for
  $\Zy\in W$. The column $u$ of $M$ is $T\Zx_u$, so the span of the columns is $T(W)$, and $T$
  is injective on $W$, since $T\Zy=0$ gives $(\Zy,\Zy)=0$. So $\dim T(W)=\dim W$ by
  \path{Submodule.map_span} and \path{LinearMap.finrank_range_of_inj}; apply
  \path{Matrix.rank_eq_finrank_span_cols}.~\qed\end{proof}
 \step{s4gfsum}{\assume{$M$ is positive semidefinite.}\prove{$M$ is the Gram matrix of a family
  $(\Zy_u)$ in $\re^m$ for some $m$.}}
 \begin{proof}\pf\ By \path{Matrix.posSemidef_iff_eq_sum_vecMulVec},
  $M=\sum_{k=1}^m\Zv_k\Zv_k^T$ for some $m$ and vectors $\Zv_k\in\re^V$. Put
  $\Zy_u=\sum_k\Zv_k(u)\Ze_k$; then $(\Zy_u,\Zy_w)=\sum_k\Zv_k(u)\Zv_k(w)$ is the $uw$-entry of
  $M$.~\qed\end{proof}
 \step{s4gfcoord}{\assume{$M$ is positive semidefinite.}\prove{$M$ is the Gram matrix of a
  family in $\re^{\dim W}$, where $W$ is the span of the $\Zy_u$ of \stepref{s4gfsum}.}}
 \begin{proof}\pf\ The coordinates in the orthonormal basis \path{stdOrthonormalBasis} of $W$,
  that is, its \path{OrthonormalBasis.repr}, form a linear isometry of $W$ onto
  $\re^{\dim W}$; they preserve inner products.~\qed\end{proof}
 \step{s4gfpad}{\assume{$k\le r$.}\prove{Padding with $r-k$ zero coordinates maps $\re^k$
  into $\re^r$ and preserves inner products.}}
 \begin{proof}\pf\ The padded coordinates contribute $0$ to each inner
  product.~\qed\end{proof}
 \qedstep
 \begin{proof}\pf\ Assertion 1 is \stepref{s4gfspan}. For assertion 2, $\dim W=\rk M\le r$ by
  \stepref{s4gfspan}; apply \stepref{s4gfcoord} and then \stepref{s4gfpad} with
  $k=\dim W$.~\qed\end{proof}
\end{proof}

\begin{namedresult}{Lemma}{(the Gram matrix)}\label{lem:4-gram}
 \lean{SimpleGraph.exists_representation_iff_posSemidef}
 \uses{def:represented,lem:4-gramfactor,facts:sec4}
 \assume{$X$ is a graph on $n$ vertices.}
 \prove{\begin{pfenum}
  \item $X$ has a representation in $\re^n$ if and only if $2I+A(X)$ is positive semidefinite,
   that is, if and only if every eigenvalue of $X$ is at least $-2$.
  \item If $n\ge1$ and $X$ has a representation in a space of dimension less than $n$, then
   $\lambda_{\min}(X)=-2$.
 \end{pfenum}}
\end{namedresult}
\begin{proof}
 \pfsketch\ We will factor the matrix into a Gram matrix, and then count independent vectors.
 \step{s4gramfwd}{\assume{$(\Zx_u)$ is a representation of $X$ in a real inner product
  space $V$.}
  \prove{$2I+A(X)$ is positive semidefinite.}}
 \begin{proof}\pf\ $2I+A(X)$ is the Gram matrix of $(\Zx_u)$ by \asm{1} and
  Definition~\ref{def:represented}; apply \path{Matrix.posSemidef_gram}.~\qed\end{proof}
 \step{s4gramback}{\assume{$2I+A(X)$ is positive semidefinite.}
  \prove{$X$ has a representation in $\re^n$.}}
 \begin{proof}\pf\ $2I+A(X)$ has rank at most $n$; apply Lemma~\ref{lem:4-gramfactor},
  assertion 2, with $r=n$, and Definition~\ref{def:represented}.~\qed\end{proof}
 \step{s4gramone}{Assertion 1 holds.}
 \begin{proof}\pf\ \stepref{s4gramfwd} and \stepref{s4gramback} give the first equivalence, and
  \path{SimpleGraph.posSemidef_two_add_adjMatrix_iff} the second.~\qed\end{proof}
 \step{s4gramsing}{\assume{$n\ge1$ and $(\Zx_u)$ is a representation of $X$ in a space $V$ of
  dimension less than $n$.}\prove{$\det(2I+A(X))=0$.}}
 \begin{proof}\pf\ The family $(\Zx_u)$ has $n$ members in a space of dimension less than $n$,
  so it is not linearly independent by \path{LinearIndependent.fintype_card_le_finrank}. Its
  Gram matrix $2I+A(X)$ is singular by
  \path{Matrix.det_gram_ne_zero_iff_linearIndependent}.~\qed\end{proof}
 \qedstep
 \begin{proof}\pf\ Assertion 1 is \stepref{s4gramone}. For assertion 2, $2I+A(X)$ is positive
  semidefinite by \stepref{s4gramfwd} and
  singular by \stepref{s4gramsing}; apply
  \path{SimpleGraph.leastEigenvalue_eq_neg_two_iff}.~\qed\end{proof}
\end{proof}

\begin{namedresult}{Lemma}{(the lines of a representation)}\label{lem:4-lines}
 \lean{SimpleGraph.IsRepresentation.isIndecomposable_of_connected}
 \uses{def:represented,def:line-system,def:indecomposable,lem:sec3-connected-span}
 \assume{$(\Zx_u)_{u\in V(X)}$ is a representation of a graph $X$ in a finite-dimensional real
 inner product space $V$, and $\Phi=\{\pm\Zx_u:u\in V(X)\}$.}
 \prove{\begin{pfenum}
  \item $u\mapsto\Zx_u$ is injective, and $\Zx_u\ne-\Zx_w$ for all $u,w$.
  \item $\Phi$ is a line system $S_\Zx$ in $V$, with $|V(X)|$ lines.
  \item If $X$ is connected, then $S_\Zx$ is indecomposable.
 \end{pfenum}}
\end{namedresult}
\begin{proof}
 \pfsketch\ We will compare inner products with the squared norm 2.
 \step{s4linesinj}{Assertion 1 holds.}
 \begin{proof}\pf\ For $u\ne w$, $(\Zx_u,\Zx_w)\in\{0,1\}$ by Definition~\ref{def:represented},
  while $(\Zx_u,\Zx_u)=2$ and $(\Zx_u,-\Zx_u)=-2$; so $\Zx_w\ne\Zx_u$ and $\Zx_w\ne-\Zx_u$. For
  $u=w$, $\Zx_u=-\Zx_u$ would give $\Zx_u=0$, against $(\Zx_u,\Zx_u)=2$.~\qed\end{proof}
 \step{s4linessys}{Assertion 2 holds.}
 \begin{proof}\pf\ $\Phi=-\Phi$, every member has squared norm 2, and the inner products of
  members are $\pm$ those of the $\Zx_u$, which lie in $\{0,1,2\}$; this is
  Definition~\ref{def:line-system}. The lines are the pairs $\{\Zx_u,-\Zx_u\}$, which are
  distinct by \stepref{s4linesinj}.~\qed\end{proof}
 \qedstep
 \begin{proof}\pf\ Take $B=\{\Zx_u\}$ in Lemma~\ref{lem:sec3-connected-span}: $B$ spans the span
  of $\Phi$, and distinct members of $B$ have nonzero inner product exactly when their vertices
  are adjacent, so the graph of that lemma is a copy of $X$, which is connected.~\qed\end{proof}
\end{proof}

\begin{namedresult}{Lemma}{(transport of representations)}\label{lem:4-transport}
 \lean{SimpleGraph.IsRepresentedIn.of_isometricInclusion}
 \uses{def:represented,def:line-iso}
 \assume{$S$ and $T$ are line systems, $f:\Phi(S)\to\Phi(T)$ preserves inner products (for
 instance an isomorphism, or the inclusion of $\Phi(S)$ in $\Phi(T)$ when $\Phi(S)\sbs\Phi(T)$),
 and $X$ is represented in $S$ by $(\Zx_u)$.}
 \prove{$X$ is represented in $T$ by $(f(\Zx_u))$.}
\end{namedresult}
\begin{proof}
 \pfsketch\ The Gram matrix does not change.
 \qedstep
 \begin{proof}\pf\ $(f(\Zx_u),f(\Zx_w))=(\Zx_u,\Zx_w)$ for all $u,w$ by \asm{0}, and
  $f(\Zx_u)\in\Phi(T)$; apply Definition~\ref{def:represented}.~\qed\end{proof}
\end{proof}

\begin{namedresult}{Lemma}{(subfamilies of a representation)}\label{lem:4-sub}
 \lean{SimpleGraph.IsRepresentation.induce}
 \uses{def:represented}
 \assume{$(\Zx_u)_{u\in V(X)}$ is a representation of a graph $X$ in $V$, and $U\sbs V(X)$.}
 \prove{$(\Zx_u)_{u\in U}$ is a representation in $V$ of the subgraph $X[U]$ induced on $U$.}
\end{namedresult}
\begin{proof}
 \pfsketch\ The Gram matrix of the subfamily is a principal submatrix.
 \qedstep
 \begin{proof}\pf\ Two vertices of $U$ are adjacent in $X[U]$ exactly when they are adjacent in
  $X$, so $(\Zx_u,\Zx_w)$ for $u,w\in U$ is the $uw$-entry of $2I+A(X[U])$; apply
  Definition~\ref{def:represented}.~\qed\end{proof}
\end{proof}

\begin{namedresult}{Lemma}{(a root system containing the lines)}\label{lem:4-rootsystem}
 \lean{SimpleGraph.IsRepresentation.exists_isRepresentedIn_rootType}
 \uses{def:represented,lem:4-lines,lem:4-transport,lem:sec3-extension,lem:sec3-transport,lem:sec3-spans,def:2.2,lem:2.3,thm:3.5,thm:3.6,def:3.1,def:3.3}
 \assume{$X$ is a connected graph with a representation $(\Zx_u)$ in a finite-dimensional real
 inner product space $V$, and the span of the $\Zx_u$ has dimension $k$.}
 \prove{\begin{pfenum}
  \item $X$ is represented in the model of one of $\typ A_k$, $\typ D_k$ ($k\ge4$), $\typ E_6$,
   $\typ E_7$ and $\typ E_8$, and the rank of that model is $k$.
  \item $X$ is represented in $\typ A_\ell$ for some $\ell\ge1$, in $\typ D_\ell$ for some
   $\ell\ge4$, or in $\typ E_8$.
 \end{pfenum}}
\end{namedresult}
\begin{proof}
 \pfsketch\ This is the paper's ``Such a set of lines is contained in a root system. If the
 graph is connected, then the root system is indecomposable.'' \repair\ The paper does not say
 why a set that is not maximal lies in an indecomposable star-closed set of the same dimension;
 we will use the fill of Gap 13 (errata), which is Lemma~\ref{lem:sec3-extension}. We will then
 place $\typ E_6$ and $\typ E_7$ inside $\typ E_8$. The first of the two sentences holds without
 connectedness: the vectors of squared norm 2 in the integral span of the $\Zx_u$ form a
 star-closed line system in which $X$ is represented. We use this system for Remark 4.14; in
 Lean it is \path{SimpleGraph.IsRepresentation.latticeRoots}, with the lemmas
 \path{isStarClosed_latticeRoots} and \path{isRepresentedIn_latticeRoots} in the same
 namespace. The claim of the introduction of [CGSS], that the lines lie in one of the root
 systems, needs connectedness (errata, Implicit hypothesis 1).
 \pflet{$S=S_\Zx$ be the line system of Lemma~\ref{lem:4-lines} and $W$ the span of $\Phi(S)$,
 which has dimension $k$.}
 \step{s4rsnonempty}{$S$ is nonempty and indecomposable.}
 \begin{proof}\pf\ A connected graph has a vertex, and Lemma~\ref{lem:4-lines}, assertion 3,
  applies by \asm{0}.~\qed\end{proof}
 \step{s4rsmax}{There is an indecomposable line system $T$ in $W$, maximal in $W$ and spanning
  $W$, with $\Phi(S)\sbs\Phi(T)$.}
 \begin{proof}
  \step{s4rsmaxyes}{\case{$S$ is maximal in $W$.}\prove{$T=S$ will do.}}
  \begin{proof}\pf\ $S$ is indecomposable by \stepref{s4rsnonempty}, maximal by \asm{2}, and
   spans $W$ by the definition of $W$.~\qed\end{proof}
  \step{s4rsmaxno}{\case{$S$ is not maximal in $W$.}\prove{Such a $T$ exists.}}
  \begin{proof}\pf\ Apply Lemma~\ref{lem:sec3-extension} to $S$, with
   \stepref{s4rsnonempty} and \asm{2}.~\qed\end{proof}
  \qedstep
  \begin{proof}\pf\ The two cases are exhaustive.~\qed\end{proof}
 \end{proof}
 \step{s4rstype}{$T$ is isomorphic to the model of $\typ A_\ell$ ($\ell\ge1$), $\typ D_\ell$
  ($\ell\ge4$), $\typ E_6$, $\typ E_7$ or $\typ E_8$, of rank $k$.}
 \begin{proof}\pf\ Apply Theorem~\ref{thm:3.6} to $T$ of \stepref{s4rsmax}; its list is part of
  this one. The isomorphism preserves rank by Lemma~\ref{lem:sec3-transport}, and $T$ has rank
  $k$ since it spans $W$.~\qed\end{proof}
 \step{s4rsinT}{$X$ is represented in $T$ by $(\Zx_u)$.}
 \begin{proof}\pf\ Lemma~\ref{lem:4-transport}, with the inclusion $\Phi(S)\sbs\Phi(T)$ of
  \stepref{s4rsmax}.~\qed\end{proof}
 \step{s4rsone}{Assertion 1 holds.}
 \begin{proof}\pf\ Apply Lemma~\ref{lem:4-transport} to \stepref{s4rsinT} and the isomorphism
  of \stepref{s4rstype}. The model of $\typ A_\ell$ or $\typ D_\ell$ has rank $\ell$ by
  Lemma~\ref{lem:sec3-spans}, so $\ell=k$.~\qed\end{proof}
 \step{s4rsE}{$\Phi(\typ E_6)\sbs\Phi(\typ E_8)$ and $\Phi(\typ E_7)\sbs\Phi(\typ E_8)$.}
 \begin{proof}\pf\ Definition~\ref{def:3.3}.~\qed\end{proof}
 \qedstep
 \begin{proof}\pf\ Assertion 1 is \stepref{s4rsone}. For assertion 2, in the cases $\typ E_6$
  and $\typ E_7$ apply Lemma~\ref{lem:4-transport} once more to the inclusions of
  \stepref{s4rsE}.~\qed\end{proof}
\end{proof}

\begin{namedresult}{Lemma}{(the line graph)}\label{lem:4-linegraph}
 \lean{SimpleGraph.lineGraph_isRepresentation}
 \uses{def:represented}
 \assume{$Y$ is a graph with vertex set $\{1,\dots,n\}$.}
 \prove{The vectors $\Ze_i+\Ze_j$ of $\re^n$, for the edges $\{i,j\}$ of $Y$, form a
 representation of $L(Y)$ in $\re^n$.}
\end{namedresult}
\begin{proof}
 \pfsketch\ The inner product counts common ends.
 \step{s4lgip}{For edges $e,f$ of $Y$, $(\sum_{i\in e}\Ze_i,\sum_{j\in f}\Ze_j)=|e\cap f|$.}
 \begin{proof}\pf\ Expand; $(\Ze_i,\Ze_j)$ is 1 if $i=j$ and 0
  otherwise.~\qed\end{proof}
 \qedstep
 \begin{proof}\pf\ $|e\cap e|=2$; for $e\ne f$, $|e\cap f|$ is 1 when $e$ and $f$ share an end,
  that is, when they are adjacent in $L(Y)$, and 0 otherwise. Apply
  \stepref{s4lgip}.~\qed\end{proof}
\end{proof}


\begin{cgss}{Theorem}{4.1}\label{thm:4.1}
 \lean{SimpleGraph.isRepresentedIn_An_iff}
 \uses{def:represented,def:3.1,lem:4-linegraph}
 \assume{$X$ is a graph and $n\ge1$.}
 \prove{$X$ is represented in $\typ A_n$ if and only if $X\cong L(Y)$ for a bipartite graph $Y$
  on the vertex set $\{1,\dots,n+1\}$.}
 \new\ The paper's ``any such graph can be represented by a subset of $\typ A_n$'' needs $Y$ to
 have at most $n+1$ vertices; we will state it with that bound, as the Lean statement does.
\end{cgss}
\begin{proof}
 \pfsketch\ We will read off $Y$ from the supports of the vectors, and show that each coordinate
 has one sign.
 \step{s41fwd}{\assume{$X$ is represented in $\typ A_n$ by $(\Zx_u)$.}
  \prove{$X\cong L(Y)$ for a bipartite graph $Y$ on $\{1,\dots,n+1\}$.}}
 \begin{proof}
  \step{s41pq}{\pick\ $p(u)\ne q(u)$ in $\{1,\dots,n+1\}$ with
   $\Zx_u=\Ze_{p(u)}-\Ze_{q(u)}$ for each vertex $u$.}
  \begin{proof}\pf\ $\Zx_u\in\Phi(\typ A_n)$ by \asm{1}; apply
   Definition~\ref{def:3.1}.~\qed\end{proof}
  \step{s41sign}{For all vertices $u,w$, $p(u)\ne q(w)$.}
  \begin{proof}\pf\ If $p(u)=q(w)=i$, then $u\ne w$, and
   $(\Zx_u,\Zx_w)=(\Ze_i-\Ze_{q(u)},\Ze_{p(w)}-\Ze_i)=-1-[q(u)=p(w)]<0$, since $p(w)\ne i$ and
   $q(u)\ne i$; this contradicts Definition~\ref{def:represented}.~\qed\end{proof}
  \step{s41ip}{For $u\ne w$, $(\Zx_u,\Zx_w)=[p(u)=p(w)]+[q(u)=q(w)]$.}
  \begin{proof}\pf\ Expand $(\Ze_{p(u)}-\Ze_{q(u)},\Ze_{p(w)}-\Ze_{q(w)})$; the two cross terms
   vanish by \stepref{s41sign}.~\qed\end{proof}
  \step{s41inj}{The supports $\{p(u),q(u)\}$ of distinct vertices are distinct.}
  \begin{proof}\pf\ If $u\ne w$ have the same support, then $p(u)=p(w)$ and $q(u)=q(w)$ by
   \stepref{s41sign}, and $(\Zx_u,\Zx_w)=2$ by \stepref{s41ip}, against
   Definition~\ref{def:represented}.~\qed\end{proof}
  \pflet{$Y$ be the graph on $\{1,\dots,n+1\}$ with edges $\{p(u),q(u)\}$, and
   $P=\{p(u):u\in V(X)\}$.}
  \step{s41bip}{$Y$ is bipartite, with parts $P$ and its complement.}
  \begin{proof}\pf\ Each edge $\{p(u),q(u)\}$ has one end $p(u)$ in $P$ and the other end
   $q(u)$ outside $P$, by \stepref{s41sign}.~\qed\end{proof}
  \step{s41iso}{$u\mapsto\{p(u),q(u)\}$ is an isomorphism from $X$ to $L(Y)$.}
  \begin{proof}\pf\ It is a bijection onto the edges of $Y$ by \stepref{s41inj}. For $u\ne w$,
   the edges share an end exactly when $p(u)=p(w)$ or $q(u)=q(w)$, by \stepref{s41sign}, that is,
   when $(\Zx_u,\Zx_w)=1$ by \stepref{s41ip}, which is adjacency of $u$ and
   $w$.~\qed\end{proof}
  \qedstep
  \begin{proof}\pf\ \stepref{s41bip} and \stepref{s41iso}.~\qed\end{proof}
 \end{proof}
 \step{s41back}{\assume{$Y$ is bipartite on $\{1,\dots,n+1\}$, with parts $P$ and $Q$, and
  $X\cong L(Y)$.}\prove{$X$ is represented in $\typ A_n$.}}
 \begin{proof}
  \pflet{$\Zx_e=\Ze_p-\Ze_q$ for the edge $e=\{p,q\}$ of $Y$ with $p\in P$ and $q\in Q$.}
  \step{s41backip}{For edges $e\ne f$, $(\Zx_e,\Zx_f)$ is 1 when $e$ and $f$ share an end, and
   0 otherwise; and $(\Zx_e,\Zx_e)=2$.}
  \begin{proof}\pf\ Write $e=\{p,q\}$ and $f=\{p',q'\}$ with $p,p'\in P$. Then $p\ne q'$ and
   $q\ne p'$, since $P\cap Q=\varnothing$, and $(\Zx_e,\Zx_f)=[p=p']+[q=q']$, which is 2 for
   $e=f$ and at most 1 for $e\ne f$.~\qed\end{proof}
  \qedstep
  \begin{proof}\pf\ $\Zx_e\in\Phi(\typ A_n)$ by Definition~\ref{def:3.1}, and $(\Zx_e)$
   represents $L(Y)$ by \stepref{s41backip}; compose with the isomorphism of
   \asm{1}.~\qed\end{proof}
 \end{proof}
 \qedstep
 \begin{proof}\pf\ \stepref{s41fwd} and \stepref{s41back}.~\qed\end{proof}
\end{proof}

\begin{namedresult}{Remark}{(after Theorem 4.1)}\label{rem:4-after41}
 \lean{LineSystem.An.card_nonneg_subset_le}\lean{LineSystem.An.nonneg_subset_maximal_iff}
 \lean{LineSystem.An.exists_nonneg_subset_card_eq}
 \lean{LineSystem.An.isRepresentation_lineGraph_completeBipartite_sigmaP}
 \uses{def:represented,def:3.1,lem:4-sub}
 \assume{$n\ge1$ and $N=n+1$. A set $\Sigma\sbs\Phi(\typ A_n)$ has \emph{nonnegative inner
 products} if $(\Zx,\Zy)\ge0$ for all $\Zx,\Zy\in\Sigma$, and is \emph{maximal} if no subset of
 $\Phi(\typ A_n)$ that properly contains it has nonnegative inner products. For
 $P\sbs\{1,\dots,N\}$ let $\Sigma_P=\{\Ze_p-\Ze_q:p\in P,\ q\notin P\}$.}
 \prove{\begin{pfenum}
  \item A set $\Sigma$ with nonnegative inner products is maximal if and only if
   $\Sigma=\Sigma_P$ for some $P$ with $1\le|P|\le n$.
  \item For such $P$, with $m=|P|$, the vectors $\Ze_p-\Ze_q$, indexed by the edges $\{p,q\}$ of
   the complete bipartite graph $Y_P$ with parts $P$ and $\{1,\dots,N\}\setminus P$, represent
   $L(Y_P)\cong L(K_{m,N-m})$.
  \item Every $\Sigma$ with nonnegative inner products has $|\Sigma|\le\lfloor\frac14N^2\rfloor$,
   and $|\Sigma_P|=\lfloor\frac14N^2\rfloor$ for $|P|=\lfloor\frac12N\rfloor$.
 \end{pfenum}}
\end{namedresult}
\begin{proof}
 \pfsketch\ \new\ The paper states the remark without proof. We will show, as in the proof of
 Theorem 4.1, that each coordinate occurs with one sign in $\Sigma$, and then count.
 \pflet{$\Zx=\Ze_{p(\Zx)}-\Ze_{q(\Zx)}$ for $\Zx\in\Phi(\typ A_n)$, by
  Definition~\ref{def:3.1}.}
 \step{s4raip}{\assume{$P\sbs\{1,\dots,N\}$, $p,p'\in P$ and $q,q'\notin P$.}
  \prove{$(\Ze_p-\Ze_q,\Ze_{p'}-\Ze_{q'})=[p=p']+[q=q']$; in particular $\Sigma_P$ has
  nonnegative inner products.}}
 \begin{proof}\pf\ Expand; the cross terms vanish since $p\ne q'$ and $q\ne
  p'$.~\qed\end{proof}
 \step{s4rasign}{\assume{$\Sigma$ has nonnegative inner products.}\prove{$\Sigma\sbs\Sigma_P$
  for $P=\{p(\Zx):\Zx\in\Sigma\}$.}}
 \begin{proof}\pf\ We need $q(\Zx)\notin P$ for $\Zx\in\Sigma$. If $q(\Zx)=p(\Zy)=i$ with
  $\Zy\in\Sigma$, then $p(\Zx)\ne i$ and $q(\Zy)\ne i$, so
  $(\Zx,\Zy)=(\Ze_{p(\Zx)}-\Ze_i,\Ze_i-\Ze_{q(\Zy)})=-1-[p(\Zx)=q(\Zy)]<0$, against
  \asm{1}.~\qed\end{proof}
 \step{s4ramax}{\assume{$\Sigma$ is maximal with nonnegative inner products.}
  \prove{$\Sigma=\Sigma_P$ for some $P$ with $1\le|P|\le n$.}}
 \begin{proof}
  \step{s4ramaxne}{$\Sigma$ is nonempty.}
  \begin{proof}\pf\ $\{\Ze_1-\Ze_2\}$ has nonnegative inner products, as $N\ge2$, and
   properly contains $\varnothing$; apply \asm{1}.~\qed\end{proof}
  \step{s4ramaxeq}{$\Sigma=\Sigma_P$ for $P=\{p(\Zx):\Zx\in\Sigma\}$.}
  \begin{proof}\pf\ $\Sigma\sbs\Sigma_P$ by \stepref{s4rasign}, and $\Sigma_P$ has nonnegative
   inner products by \stepref{s4raip}; apply \asm{1}.~\qed\end{proof}
  \qedstep
  \begin{proof}\pf\ By \stepref{s4ramaxne}, $P$ of \stepref{s4ramaxeq} is nonempty, and its
   complement contains $q(\Zx)$ for $\Zx\in\Sigma=\Sigma_P$; so $1\le|P|\le
   N-1=n$.~\qed\end{proof}
 \end{proof}
 \step{s4raback}{\assume{$1\le|P|\le n$.}\prove{$\Sigma_P$ is maximal with nonnegative inner
  products.}}
 \begin{proof}\pf\ $\Sigma_P$ has nonnegative inner products by \stepref{s4raip}. A root
  $\Zz=\Ze_a-\Ze_b\notin\Sigma_P$ has $b\in P$ or $a\notin P$; pick $p\in P$ and $q\notin P$, by
  \asm{1}. If $b\in P$, then $(\Zz,\Ze_b-\Ze_q)=-1-[a=q]<0$; if $a\notin P$, then
  $(\Zz,\Ze_p-\Ze_a)=-1-[b=p]<0$. Thus no larger set has nonnegative inner
  products.~\qed\end{proof}
 \step{s4rarep}{Assertion 2 holds.}
 \begin{proof}\pf\ For edges $e=\{p,q\}$ and $f=\{p',q'\}$ of $Y_P$ with $p,p'\in P$,
  \stepref{s4raip} gives the inner product $[p=p']+[q=q']=|e\cap f|$, which is 2 for $e=f$, and
  for $e\ne f$ is 1 when $e$ and $f$ share an end and 0 otherwise; apply
  Definition~\ref{def:represented}. $Y_P\cong K_{m,N-m}$.~\qed\end{proof}
 \step{s4racard}{$|\Sigma_P|=|P|(N-|P|)$.}
 \begin{proof}\pf\ $(p,q)\mapsto\Ze_p-\Ze_q$ is injective on $P\times(\{1,\dots,N\}\setminus
  P)$.~\qed\end{proof}
 \step{s4rabound}{\assume{$\Sigma$ has nonnegative inner products.}
  \prove{$|\Sigma|\le\lfloor\frac14N^2\rfloor$.}}
 \begin{proof}
  \pflet{$P$ be as in \stepref{s4rasign}, and $m=|P|$.}
  \step{s4rabsub}{$|\Sigma|\le m(N-m)$.}
  \begin{proof}\pf\ $\Sigma\sbs\Sigma_P$ by \asm{1} and \stepref{s4rasign}, and
   $|\Sigma_P|=m(N-m)$ by \stepref{s4racard}.~\qed\end{proof}
  \step{s4rabsq}{$m(N-m)\le\lfloor\frac14N^2\rfloor$.}
  \begin{proof}\pf\ $4m(N-m)=N^2-(N-2m)^2\le N^2$, and $m(N-m)$ is an
   integer.~\qed\end{proof}
  \qedstep
  \begin{proof}\pf\ \stepref{s4rabsub} and \stepref{s4rabsq}.~\qed\end{proof}
 \end{proof}
 \step{s4raeq}{\assume{$m=\lfloor\frac12N\rfloor$.}\prove{$m(N-m)=\lfloor\frac14N^2\rfloor$.}}
 \begin{proof}\pf\ For $N=2k$ both sides are $k^2$; for $N=2k+1$ both are
  $k^2+k$.~\qed\end{proof}
 \qedstep
 \begin{proof}\pf\ Assertion 1 is \stepref{s4ramax} and \stepref{s4raback}, and assertion 2 is
  \stepref{s4rarep}. Assertion 3 is \stepref{s4rabound}, with \stepref{s4racard} and
  \stepref{s4raeq} for $|P|=\lfloor\frac12N\rfloor$, which lies between 1 and $n$ as
  $N\ge2$.~\qed\end{proof}
\end{proof}
The paper states assertions 1 to 3 in the form ``The maximal subsets of $\typ A_n$ with
nonnegative inner products represent the graphs $L(K_{m,n+1-m})$, and have maximum cardinality
$\lfloor\frac14(n+1)^2\rfloor$, attained for $m=\lfloor\frac12(n+1)\rfloor$.'' A set
$\Sigma$ with nonnegative inner products lies in $\Sigma_P$ for $P=\{p(\Zx):\Zx\in\Sigma\}$, by
the second step of the proof. Thus, by assertion 2 and Lemma~\ref{lem:4-sub}, the graph that
$\Sigma$ represents is $L(Y)$, where $Y$ is the subgraph of $Y_P$ whose edges are the
$\{p,q\}$ with $\Ze_p-\Ze_q\in\Sigma$, and assertion 3 bounds its number of vertices.

\subsection{Theorems 4.2 to 4.9}

\begin{namedresult}{Definition}{(cocktail party graphs)}\label{def:CP}
 \lean{SimpleGraph.cocktailParty}
 \uses{}
 \assume{$m\ge0$ is an integer.}
 \pflet{The \emph{cocktail party graph} $CP(m)$ has vertex set $\{1,\dots,m\}\times\{0,1\}$, and
 $(i,s)$ is adjacent to $(j,t)$ exactly when $i\ne j$.}
 Thus $CP(m)$ is $K_{2m}$ with a perfect matching deleted, and it is $(2m-2)$-regular on $2m$
 vertices; $CP(0)$ has no vertices.
\end{namedresult}

\begin{namedresult}{Definition}{(generalized line graphs)}\label{def:glg}
 \lean{SimpleGraph.generalizedLineGraph}\lean{SimpleGraph.IsGeneralizedLineGraph}
 \uses{def:CP}
 \assume{$Y$ is a graph with vertex set $\{1,\dots,m\}$ and $a_1,\dots,a_m\ge0$ are integers.}
 \pflet{The \emph{generalized line graph} $L(Y;a_1,\dots,a_m)$ has as vertices the edges of $Y$
 and the triples $(i,\alpha,s)$ with $1\le i\le m$, $1\le\alpha\le a_i$ and $s\in\{0,1\}$. Two
 edges are adjacent when they are adjacent in $L(Y)$; $(i,\alpha,s)$ and $(j,\beta,t)$ are
 adjacent when $i=j$ and $\alpha\ne\beta$; and $(i,\alpha,s)$ is adjacent to the edge $e$ when
 $i\in e$. A graph is a \emph{generalized line graph} if it is isomorphic to some
 $L(Y;a_1,\dots,a_m)$.}
 For each $i$, the triples $(i,\alpha,s)$ span a copy of $CP(a_i)$, and $L(Y;0,\dots,0)\cong
 L(Y)$.
\end{namedresult}

\begin{namedresult}{Lemma}{(the representation of a generalized line graph)}\label{lem:4-glgrep}
 \lean{SimpleGraph.generalizedLineGraph_isRepresentedIn_Dn}
 \uses{def:glg,def:represented,def:3.1,lem:4-linegraph}
 \assume{$Y$ and $a_1,\dots,a_m$ are as in Definition~\ref{def:glg}, $n=\sum_{i=1}^m(1+a_i)$,
 and the orthonormal basis of $\re^n$ is written $\Ze_{i,\alpha}$, for $1\le i\le m$ and
 $0\le\alpha\le a_i$.}
 \prove{The vectors $\Zx_e=\Ze_{i,0}+\Ze_{j,0}$ for the edges $e=\{i,j\}$ of $Y$, and
 $\Zx_{(i,\alpha,s)}=\Ze_{i,0}+(-1)^s\Ze_{i,\alpha}$, lie in $\Phi(\typ D_n)$ and represent
 $L(Y;a_1,\dots,a_m)$.}
\end{namedresult}
\begin{proof}
 \pfsketch\ \repair\ We read the paper's ``$1\le\alpha\le\alpha_i$'' as $1\le\alpha\le a_i$
 (errata, Error 13, item 5). We will compute the three kinds of inner products.
 \step{s4glgee}{For edges $e,f$, $(\Zx_e,\Zx_f)=|e\cap f|$.}
 \begin{proof}\pf\ Lemma~\ref{lem:4-linegraph}, in the coordinates $\Ze_{i,0}$, gives 2 for
  $e=f$, 1 when $e\ne f$ share an end, and 0 otherwise; this is $|e\cap f|$.~\qed\end{proof}
 \step{s4glgce}{$(\Zx_{(i,\alpha,s)},\Zx_e)$ is 1 if $i\in e$ and 0 otherwise.}
 \begin{proof}\pf\ $\Ze_{i,\alpha}$ with $\alpha\ge1$ does not occur in $\Zx_e$, and
  $(\Ze_{i,0},\Zx_e)=[i\in e]$.~\qed\end{proof}
 \step{s4glgcc}{$(\Zx_{(i,\alpha,s)},\Zx_{(j,\beta,t)})$ is 2 if the triples are equal, 0 if
  $i=j$ and $\alpha=\beta$ and $s\ne t$, 1 if $i=j$ and $\alpha\ne\beta$, and 0 if $i\ne j$.}
 \begin{proof}\pf\ The inner product is
  $[i=j]\bigl(1+(-1)^{s+t}[\alpha=\beta]\bigr)$.~\qed\end{proof}
 \qedstep
 \begin{proof}\pf\ Each vector is a sum or difference of two distinct basis vectors, so lies in
  $\Phi(\typ D_n)$ by Definition~\ref{def:3.1}. By \stepref{s4glgee}, \stepref{s4glgce} and
  \stepref{s4glgcc}, the Gram matrix is $2I$ plus the adjacency matrix of
  Definition~\ref{def:glg}.~\qed\end{proof}
\end{proof}


\begin{namedresult}{Lemma}{(the normal form in $\typ D_n$)}\label{lem:4-Dnform}
 \lean{SimpleGraph.IsRepresentation.exists_normal_form_Dn}
 \uses{def:represented,def:3.1,def:glg,lem:4-glgrep,lem:4-transport}
 \assume{$n\ge1$, and $(\Zx_u)_{u\in V(X)}$ represents a graph $X$ in $\typ D_n$.}
 \prove{There are a set $C\sbs\{1,\dots,n\}$, integers $a_c\ge0$ for $c\in C$, distinct
 coordinates $p_{c,\alpha}$ for $c\in C$ and $1\le\alpha\le a_c$, which are the members of
 $\{1,\dots,n\}\setminus C$, and signs $\epsilon_1,\dots,\epsilon_n\in\{1,-1\}$ such that, with
 $D$ the diagonal isometry $\Ze_i\mapsto\epsilon_i\Ze_i$ and $\Zy_u=D\Zx_u$:
 \begin{pfenum}
  \item $(\Zy_u)$ represents $X$ in $\typ D_n$;
  \item for each $c\in C$, $1\le\alpha\le a_c$ and $s\in\{0,1\}$ exactly one vertex $u$ has
   $\Zy_u=\Ze_c+(-1)^s\Ze_{p_{c,\alpha}}$, and every vertex $u$ not of this form has
   $\Zy_u=\Ze_i+\Ze_j$ with $i\ne j$ in $C$, distinct such vertices having distinct supports;
  \item $X\cong L(Y;(a_c)_{c\in C})$, where $Y$ is the graph on $C$ whose edges are the supports
   of the vertices of assertion 2 not of this form.
 \end{pfenum}
 In particular $|C|+\sum_{c\in C}a_c=n$.}
 \new\ This is the proof of assertion 1 of Theorem 4.2 below, with the coordinates kept; the
 remark after that theorem counts them.
\end{namedresult}
\begin{proof}
 \pfsketch\ We will follow the paper: we will find the pairs $\Ze_c\pm\Ze_p$, then choose the
 signs of the other coordinates. For $n<4$, $\typ D_n$ is the coordinate model
 $\{\pm\Ze_i\pm\Ze_j:i\ne j\}$, as in the remark of Definition~\ref{def:3.1}.
 \repair\ The paper leaves both steps as assertions, among them that the set of pairs is
 unique and maximal (errata, Gap 19); we will prove what we need.
 \step{s42form}{Each $\Zx_u$ is $\sigma\Ze_i+\tau\Ze_j$ with $i\ne j$ and
  $\sigma,\tau\in\{1,-1\}$; its support has two members.}
 \begin{proof}\pf\ \asm{0} and Definition~\ref{def:3.1}.~\qed\end{proof}
 \step{s42same}{\assume{$u\ne w$ and $\supp(\Zx_u)=\supp(\Zx_w)=\{i,j\}$.}
  \prove{$(\Zx_u,\Zx_w)=0$, and $\Zx_u,\Zx_w$ agree in exactly one coordinate.}}
 \begin{proof}\pf\ The inner product is $\sigma\sigma'+\tau\tau'\in\{-2,0,2\}$ by
  \stepref{s42form}, and lies in $\{0,1\}$ by \asm{0} and Definition~\ref{def:represented}; so it is 0, and
  exactly one of $\sigma\sigma'$, $\tau\tau'$ is 1.~\qed\end{proof}
 \pflet{A \emph{pair} is a set $\{u,w\}$ as in \stepref{s42same}. Its \emph{centre} is the
  coordinate in which $\Zx_u$ and $\Zx_w$ agree and its \emph{petal} the other one; let $\Pi$
  be the set of petals of all pairs.}
 \step{s42petal}{\assume{$p\in\Pi$ is the petal of the pair $\{u,w\}$, with centre $c$.}
  \prove{$u$ and $w$ are the only vertices $z$ with $p\in\supp(\Zx_z)$.}}
 \begin{proof}
  \step{s42petalsupp}{\assume{$p\in\supp(\Zx_z)$.}\prove{$\supp(\Zx_z)=\{c,p\}$.}}
  \begin{proof}\pf\ Otherwise $\Zx_z=\rho\Ze_p+\ldots$ has no $\Ze_c$ term, so
   $(\Zx_z,\Zx_u)=\rho\tau$ and $(\Zx_z,\Zx_w)=-\rho\tau$, where $\pm\tau$ are the coefficients
   of $\Ze_p$ in $\Zx_u,\Zx_w$; one of them is $-1$, against \asm{0} and
   Definition~\ref{def:represented}.~\qed\end{proof}
  \qedstep
  \begin{proof}\pf\ If $z\ne u,w$, then $z$, $u$ and $w$ have pairwise the support $\{c,p\}$ by
   \stepref{s42petalsupp}, so pairwise agree in exactly one of two coordinates by
   \stepref{s42same}. Their three sign vectors in $\{1,-1\}^2$ would then be pairwise adjacent in
   the graph on $\{1,-1\}^2$ joining vectors that differ in exactly one coordinate; that graph
   is $C_4$, which has no triangle.~\qed\end{proof}
 \end{proof}
 \step{s42centre}{No centre lies in $\Pi$, and every petal is the petal of exactly one
  pair.}
 \begin{proof}\pf\ The pair of a petal is the set of vertices whose support contains it, by
  \stepref{s42petal}. If the centre $c$ of the pair $\{u,w\}$ were the petal of a pair, that
  pair would be $\{u,w\}$ by the same step, whose petal is not $c$.~\qed\end{proof}
 \pflet{$C=\{1,\dots,n\}\setminus\Pi$.}
 \step{s42onesign}{\assume{$i\in C$.}\prove{All $\Zx_u$ with $i\in\supp(\Zx_u)$ have the same
  coefficient $\epsilon_i$ at $\Ze_i$.}}
 \begin{proof}\pf\ Suppose $\Zx_u$ and $\Zx_w$ have coefficients $1$ and $-1$ at $\Ze_i$. If
  their supports are equal, $\{u,w\}$ is a pair that disagrees at $i$, so $i\in\Pi$ by the
  definition of petals; otherwise $(\Zx_u,\Zx_w)=-1$. The first contradicts \asm{1}, and the second \asm{0}
  and Definition~\ref{def:represented}.~\qed\end{proof}
 \pflet{$\epsilon_i=1$ for the $i\in C$ in no support, and $\epsilon_p=1$ for $p\in\Pi$; let
  $D$ be the diagonal isometry $\Ze_i\mapsto\epsilon_i\Ze_i$, and $\Zy_u=D\Zx_u$.}
 \step{s42flip}{$(\Zy_u)$ represents $X$ in $\typ D_n$; every $\Zy_u$ is $\Ze_c\pm\Ze_p$ for a
  pair with centre $c$ and petal $p$, or $\Ze_i+\Ze_j$ with $i,j\in C$.}
 \begin{proof}\pf\ $D$ is an isometry that preserves $\Phi(\typ D_n)$, so
  Lemma~\ref{lem:4-transport} applies. A vertex in a pair has support $\{c,p\}$, and
  $c\in C$ by \stepref{s42centre}, so its $\Ze_c$ coefficient becomes 1 by
  \stepref{s42onesign}. Every other support lies in $C$ by \stepref{s42petal}, and both
  coefficients become 1 by \stepref{s42onesign}.~\qed\end{proof}
 \pflet{$a_c$ be the number of pairs with centre $c$, for $c\in C$; number the petals of centre
  $c$ as $p_{c,1},\dots,p_{c,a_c}$; and let $Y$ be the graph on $C$ whose edges are the supports
  of the vertices in no pair.}
 \step{s42petals}{The $p_{c,\alpha}$ are distinct, and they are the members of
  $\{1,\dots,n\}\setminus C$.}
 \begin{proof}\pf\ Each petal has one pair by \stepref{s42centre}, and so one centre, which lies
  in $C$ by the same step; and $\{1,\dots,n\}\setminus C=\Pi$.~\qed\end{proof}
 \step{s42two}{Assertion 2 holds.}
 \begin{proof}\pf\ The vertices of the pair with petal $p_{c,\alpha}$ have
  $\Zy_u=\Ze_c\pm\Ze_{p_{c,\alpha}}$ by \stepref{s42flip}, one with each sign by
  \stepref{s42same}, and no other vertex has $p_{c,\alpha}$ in its support by
  \stepref{s42petal}. The vertices in no pair have $\Zy_u=\Ze_i+\Ze_j$ with $i,j\in C$ by
  \stepref{s42flip}, and distinct ones have distinct supports, since equal supports make a
  pair.~\qed\end{proof}
 \step{s42bij}{The map $\psi$ sending a vertex in no pair to its support, and the vertex $u$
  with $\Zy_u=\Ze_c+(-1)^s\Ze_{p_{c,\alpha}}$ to $(c,\alpha,s)$, is a bijection from $V(X)$
  to the vertex set of $L(Y;(a_c)_{c\in C})$.}
 \begin{proof}\pf\ By \stepref{s42two}, every triple $(c,\alpha,s)$ is the image of exactly one
  vertex, and the other vertices map injectively onto the edges of $Y$, by the definition of
  $Y$.~\qed\end{proof}
 \step{s42three}{Assertion 3 holds.}
 \begin{proof}\pf\ Renumber $C$ as $\{1,\dots,m\}$. The vectors $\Zy_u$ are those of
  Lemma~\ref{lem:4-glgrep} for $\psi(u)$, with $\Ze_{c,0}=\Ze_c$ and
  $\Ze_{c,\alpha}=\Ze_{p_{c,\alpha}}$, by \stepref{s42two} and \stepref{s42petals}. Both
  families have Gram matrix $2I$ plus the adjacency matrix, of $X$ by \stepref{s42flip} and of
  $L(Y;a)$ by that lemma; so $\psi$ of \stepref{s42bij} preserves adjacency.~\qed\end{proof}
 \qedstep
 \begin{proof}\pf\ Assertion 1 is part of \stepref{s42flip}, and assertions 2 and 3 are
  \stepref{s42two} and \stepref{s42three}. The last sentence counts $\{1,\dots,n\}$ by
  \stepref{s42petals}.~\qed\end{proof}
\end{proof}


\begin{cgss}{Theorem}{4.2}\label{thm:4.2}
 \lean{SimpleGraph.isRepresentedIn_Dn_iff}
 \uses{def:represented,def:3.1,def:glg,lem:4-glgrep,lem:4-Dnform}
 \assume{$X$ is a graph and $n\ge1$.}
 \prove{\begin{pfenum}
  \item If $X$ is represented in $\typ D_n$, then $X$ is a generalized line graph.
  \item $L(Y;a_1,\dots,a_m)$ is represented in $\typ D_n$ for $n=\sum_i(1+a_i)$.
 \end{pfenum}}
\end{cgss}
\begin{proof}
 \pfsketch\ Assertion 2 is Lemma~\ref{lem:4-glgrep}. For assertion 1 we will follow the paper:
 we will find the pairs $\Ze_c\pm\Ze_p$, then choose the signs of the other coordinates; this is
 Lemma~\ref{lem:4-Dnform}, which also fills Gap 19 (errata).
 \repair\ The paper's ``whose lines from a subset'' is ``form'' (errata, Error 13, item 4).
 \step{s42glg}{Assertion 2 holds.}
 \begin{proof}\pf\ Lemma~\ref{lem:4-glgrep}.~\qed\end{proof}
 \step{s42one}{\assume{$(\Zx_u)$ represents $X$ in $\typ D_n$.}
  \prove{$X$ is a generalized line graph.}}
 \begin{proof}\pf\ Lemma~\ref{lem:4-Dnform}, assertion 3, with \asm{1}, and $n\ge1$ by
  \asm{0}; then Definition~\ref{def:glg}.~\qed\end{proof}
 \qedstep
 \begin{proof}\pf\ \stepref{s42one} and \stepref{s42glg}.~\qed\end{proof}
\end{proof}

\begin{namedresult}{Remark}{(after Theorem 4.2)}\label{rem:4-after42}
 \lean{LineSystem.Dn.card_nonneg_subset_le}
 \lean{LineSystem.Dn.exists_isRepresentation_generalizedLineGraph_top_of_maximal}
 \lean{LineSystem.Dn.exists_nonneg_subset_card_eq}
 \uses{def:represented,def:3.1,def:glg,lem:4-Dnform,lem:4-linegraph,rem:4-after41}
 \assume{$n\ge1$, and $\Sigma\sbs\Phi(\typ D_n)$ has nonnegative inner products, in the sense of
 \hyperref[rem:4-after41]{the remark after Theorem 4.1}, with $\typ D_n$ for $\typ A_n$.}
 \prove{\begin{pfenum}
  \item If $\Sigma$ is maximal, then there are $m\ge1$ and $a_1,\dots,a_m\ge0$ with
   $\sum_{i=1}^m(1+a_i)=n$ such that the vectors of $\Sigma$, suitably indexed, represent
   $L(K_m;a_1,\dots,a_m)$.
  \item If $n\ge4$, then $|\Sigma|\le\frac12n(n-1)$; the set $\{\Ze_i+\Ze_j:i<j\}$ has
   nonnegative inner products and $\frac12n(n-1)$ members, and represents $L(K_n)$.
 \end{pfenum}}
\end{namedresult}
\begin{proof}
 \pfsketch\ \new\ The paper states the remark without proof. We will apply
 Lemma~\ref{lem:4-Dnform}, which is the proof of Theorem 4.2 with the coordinates
 kept, and count. \new\ The bound fails for $n=2$ and $n=3$, where the paper does not define
 $\typ D_n$: the sets $\{\Ze_1\pm\Ze_2\}$ and $\{\Ze_1\pm\Ze_2,\Ze_1\pm\Ze_3\}$ are larger than
 $\frac12n(n-1)$.
 \step{s4rdrep}{The family $(\Zx)_{\Zx\in\Sigma}$ represents in $\typ D_n$ the graph $X$ on
  $\Sigma$ in which $\Zx\sim\Zy$ when $(\Zx,\Zy)=1$.}
 \begin{proof}\pf\ For $\Zx\ne\Zy$ in $\Sigma$, $(\Zx,\Zy)$ is an integer at least 0 by
  \asm{0} and Definition~\ref{def:3.1}, and $(\Zx,\Zy)<2$ since
  $(\Zx-\Zy,\Zx-\Zy)=4-2(\Zx,\Zy)>0$; so it is 0 or 1. Apply
  Definition~\ref{def:represented}.~\qed\end{proof}
 \step{s4rdnf}{\pick\ $C$, $(a_c)$, $(p_{c,\alpha})$, $D$, $(\Zy_\Zx)$ and $Y$ as in
  Lemma~\ref{lem:4-Dnform} for the representation of \stepref{s4rdrep}, and let $m=|C|$. Then
  $m+\sum_ca_c=n$ and $m\ge1$.}
 \begin{proof}\pf\ The equation is the last sentence of the lemma. If $m=0$, then
  $C=\varnothing$, so the equation gives $n=0$, against \asm{0}.~\qed\end{proof}
 \step{s4rdcount}{$|\Sigma|=|E(Y)|+2\sum_ca_c\le\frac12m(m-1)+2(n-m)$.}
 \begin{proof}\pf\ Lemma~\ref{lem:4-Dnform}, assertion 2, gives two members of $\Sigma$ for each
  $p_{c,\alpha}$ and one for each edge of $Y$, a graph on the $m$ vertices of $C$; and $\sum_c
  a_c=n-m$ by \stepref{s4rdnf}.~\qed\end{proof}
 \step{s4rdmax}{\assume{$\Sigma$ is maximal.}\prove{$Y$ is complete.}}
 \begin{proof}
  \step{s4rdadd}{\assume{$i\ne j$ lie in $C$.}\prove{$\Zz=D(\Ze_i+\Ze_j)$ lies in
   $\Phi(\typ D_n)$, and $(\Zz,\Zx)\ge0$ for every $\Zx\in\Sigma$.}}
  \begin{proof}
   \step{s4rdaddroot}{$\Zz=\epsilon_i\Ze_i+\epsilon_j\Ze_j$, and $\Zz\in\Phi(\typ D_n)$.}
   \begin{proof}\pf\ $D$ is the diagonal isometry $\Ze_k\mapsto\epsilon_k\Ze_k$, and $i\ne j$ by
    \asm{2}; apply Definition~\ref{def:3.1}.~\qed\end{proof}
   \step{s4rdaddip}{\assume{$\Zx\in\Sigma$.}\prove{$(\Zz,\Zx)=(\Ze_i+\Ze_j,\Zy_\Zx)$.}}
   \begin{proof}\pf\ $D$ is an isometry with $D^2=I$, so
    $(D(\Ze_i+\Ze_j),\Zx)=(\Ze_i+\Ze_j,D\Zx)$, and $D\Zx=\Zy_\Zx$.~\qed\end{proof}
   \step{s4rdaddnn}{\assume{$\Zx\in\Sigma$.}\prove{$(\Ze_i+\Ze_j,\Zy_\Zx)\ge0$.}}
   \begin{proof}\pf\ By Lemma~\ref{lem:4-Dnform}, assertion 2, $\Zy_\Zx$ is
    $\Ze_c\pm\Ze_{p_{c,\alpha}}$ with $p_{c,\alpha}\notin C$, which gives $[c\in\{i,j\}]$ by
    \asm{2}, or $\Ze_k+\Ze_l$, which gives $|\{i,j\}\cap\{k,l\}|$.~\qed\end{proof}
   \qedstep
   \begin{proof}\pf\ \stepref{s4rdaddroot}, \stepref{s4rdaddip} and
    \stepref{s4rdaddnn}.~\qed\end{proof}
  \end{proof}
  \qedstep
  \begin{proof}\pf\ For $i\ne j$ in $C$, $\Zz\in\Sigma$ by \stepref{s4rdadd} and \asm{1}. Its
   image $\Zy_\Zz=\Ze_i+\Ze_j$ has no petal $p_{c,\alpha}$ in its support, so it is of the second
   kind of Lemma~\ref{lem:4-Dnform}, assertion 2, and $\{i,j\}$ is an edge of
   $Y$.~\qed\end{proof}
 \end{proof}
 \step{s4rdone}{Assertion 1 holds.}
 \begin{proof}\pf\ Renumber $C$ as $\{1,\dots,m\}$. By Lemma~\ref{lem:4-Dnform}, assertion 3,
  and \stepref{s4rdmax}, $X\cong L(K_m;a_1,\dots,a_m)$, and $\sum_i(1+a_i)=n$ by
  \stepref{s4rdnf}; index $\Sigma$ by this isomorphism.~\qed\end{proof}
 \step{s4rdineq}{\assume{$n\ge4$ and $1\le m\le n$.}
  \prove{$\frac12m(m-1)+2(n-m)\le\frac12n(n-1)$.}}
 \begin{proof}\pf\ The difference of the two sides is $\frac12(n-m)(n+m-5)$, and $n+m\ge5$ by
  \asm{1}.~\qed\end{proof}
 \step{s4rdK}{$\{\Ze_i+\Ze_j:i<j\}$ has $\frac12n(n-1)$ members and nonnegative inner products,
  and represents $L(K_n)$.}
 \begin{proof}\pf\ Lemma~\ref{lem:4-linegraph} for $Y=K_n$, whose Gram matrix $2I+A$ has no
  negative entry; the vectors lie in $\Phi(\typ D_n)$ by
  Definition~\ref{def:3.1}.~\qed\end{proof}
 \qedstep
 \begin{proof}\pf\ Assertion 1 is \stepref{s4rdone}. For assertion 2, \stepref{s4rdcount} and
  \stepref{s4rdineq}, with $1\le m\le n$ by \stepref{s4rdnf}, give the bound, and
  \stepref{s4rdK} the example.~\qed\end{proof}
\end{proof}
The paper states the remark as ``A maximal subset of $\typ D_n$ with nonnegative inner products
represents $L(K_m;\mathbf a)$ for some $m$ and $\mathbf a$; the maximum cardinality is
$\frac12n(n-1)$, attained by $L(K_n)$.'' Assertion 1 holds for every $n\ge1$. For $n=4$ the
maximum 6 is attained also by sets with $(m;a_1)=(1;3)$; these represent $L(K_1;3)=CP(3)$,
which is isomorphic to $L(K_4)$, so what differs is the set and its parameters, not the graph.
For $n\ge5$ the maximum is attained only by sets with $m=n$, which represent $L(K_n)$, since
$\frac12(n-m)(n+m-5)>0$ for $1\le m<n$.


\begin{namedresult}{Lemma}{(the symmetric square)}\label{lem:4-symsq}
 \lean{EuclideanSpace.inner_symSq}
 \uses{}
 \assume{$\ell\ge1$, and for $\Zv\in\re^\ell$ the vector $\sigma(\Zv)\in\re^{\ell(\ell+1)/2}$
 has the coordinates $\Zv(i)^2$ for $1\le i\le\ell$ and $\sqrt2\,\Zv(i)\Zv(j)$ for\\ $1\le i<j\le\ell$.}
 \prove{$(\sigma(\Zv),\sigma(\Zw))=(\Zv,\Zw)^2$ for all $\Zv,\Zw\in\re^\ell$.}
 \new\ The paper uses the projections onto the lines, in the space of symmetric maps with the
 trace inner product; the coordinates of that space give $\sigma$, and $\re^{36}$ is easier in
 Lean.
\end{namedresult}
\begin{proof}
 \pfsketch\ We will expand the square.
 \qedstep
 \begin{proof}\pf\ $(\Zv,\Zw)^2=\bigl(\sum_i\Zv(i)\Zw(i)\bigr)^2=\sum_i\Zv(i)^2\Zw(i)^2+2\sum_{i<j}\Zv(i)\Zv(j)\Zw(i)\Zw(j)$, which
  is $(\sigma(\Zv),\sigma(\Zw))$.~\qed\end{proof}
\end{proof}

\begin{namedresult}{Lemma}{(valencies from the sets $N_\Zr$)}\label{lem:4-valency}
 \lean{SimpleGraph.IsRepresentedIn.degree_le_of_card_N_le}
 \uses{def:2.1,def:represented,lem:4-lines,def:line-system,facts:sec4}
 \assume{\begin{pfenum}
  \item $S$ is a star-closed line system, $b\ge0$, and $|N_\Zr|\le2b$ for every
   $\Zr\in\Phi(S)$, where $N_\Zr=\{\Zx\in\Phi(S):(\Zx,\Zr)=1\}$.
  \item $X$ is represented in $S$ by $(\Zx_u)$, and $u$ is a vertex of $X$.
 \end{pfenum}}
 \prove{$u$ has at most $b$ neighbours.}
 \repair\ The paper cites ``Lemma 3.7'' for such bounds (errata, Error 6), and Theorem 3.7 gives
 only the number $|N_\Zr|$; the pairing below is the fill of Gap 14 (errata), with LN6 of
 \path{cgss/prerequisites.md}.
\end{namedresult}
\begin{proof}
 \pfsketch\ We will pair the vectors with inner product 1 with $\Zr=\Zx_u$.
 \pflet{$\Zr=\Zx_u$.}
 \step{s4valpair}{\assume{$\Zx\in N_\Zr$.}\prove{$\Zr-\Zx\in N_\Zr$, $\Zr-\Zx\ne\Zx$ and
  $(\Zx,\Zr-\Zx)=-1$.}}
 \begin{proof}\pf\ Since $(\Zx,-\Zr)=-1$ and $S$ is star-closed by \asm[1]{0}, the vector
  $\Zx-\Zr$ and its negative lie in $\Phi(S)$. Further $(\Zr-\Zx,\Zr)=2-1=1$ and
  $(\Zx,\Zr-\Zx)=1-2=-1$, which also gives $\Zr-\Zx\ne\Zx$.~\qed\end{proof}
 \step{s4valnbr}{$w\mapsto\Zx_w$ maps the neighbours of $u$ injectively into $N_\Zr$, and its
  image meets each pair $\{\Zx,\Zr-\Zx\}$ at most once.}
 \begin{proof}\pf\ $(\Zx_w,\Zr)=1$ for a neighbour $w$ by Definition~\ref{def:represented}, and
  the map is injective by Lemma~\ref{lem:4-lines}, assertion 1. Two neighbours have inner product
  0 or 1, not the $-1$ of \stepref{s4valpair}.~\qed\end{proof}
 \qedstep
 \begin{proof}\pf\ By \stepref{s4valpair}, $\Zx\mapsto\Zr-\Zx$ is an involution of $N_\Zr$
  without fixed points, so $N_\Zr$ is the union of at most $b$ disjoint pairs by
  \asm[1]{0}; apply \stepref{s4valnbr}.~\qed\end{proof}
\end{proof}

\begin{namedresult}{Computation}{C20}\label{comp:C20}
 \lean{LineSystem.An.card_N}\lean{LineSystem.Dn.card_N}\lean{LineSystem.E₆.card_N}
 \lean{LineSystem.E₇.card_N}\lean{LineSystem.E₈.card_N}
 \uses{def:3.1,def:3.2,def:3.3,thm:3.5,lem:4-valency}
 \assume{$\Zr$ is a vector of the model of $\typ A_\ell$, $\typ D_\ell$, $\typ E_6$,
 $\typ E_7$ or $\typ E_8$, and $X$ is represented in that model.}
 \prove{$|N_\Zr|$ is $2(\ell-1)$, $4(\ell-2)$, 20, 32 or 56 respectively, and every vertex of
 $X$ has at most $\ell-1$, $2(\ell-2)$, 10, 16 or 28 neighbours respectively.}
\end{namedresult}
\begin{proof}
 \pfsketch\ The counts are LN6 of \path{cgss/prerequisites.md}. We will cite their Lean
 proofs, which are parametric in $\ell$ for $\typ A_\ell$ and $\typ D_\ell$ and use
 \texttt{decide} for the three exceptional models.
 \step{s4c20count}{The counts of $|N_\Zr|$ hold.}
 \begin{proof}\pf\ \path{LineSystem.An.card_N}, \path{LineSystem.Dn.card_N}, and
  \leancite{LineSystem.E₆.card_N}{\texttt{LineSystem.E\ensuremath{{}_6}.card\_N}} and its
  analogues for $\typ E_7$ and $\typ E_8$.~\qed\end{proof}
 \qedstep
 \begin{proof}\pf\ The models are star-closed by Theorem~\ref{thm:3.5}; apply
  Lemma~\ref{lem:4-valency} with half of each count.~\qed\end{proof}
\end{proof}

\begin{namedresult}{Computation}{C7}\label{comp:C7}
 \lean{Computations.C7.isRepresentedIn_E₈}
 \uses{def:3.1,def:3.2,def:represented}
 \assume{In $\re^9$, let $\Zc=\frac13\sum_{l=1}^9\Ze_l$, $\Zy_h=\Ze_h-\Ze_9$ for $1\le h\le8$,
 and $\Zz_{ij}=\Zc-\Ze_i-\Ze_j-\Ze_9$ for $1\le i<j\le8$.}
 \prove{These 36 vectors lie in model (i) of $\typ E_8$ and represent a graph $X_{36}$ on 36
 vertices in which $\Zy_1$ has 28 neighbours.}
\end{namedresult}
\begin{proof}
 \pfsketch\ We will compute the inner products; the Lean proof is by \texttt{decide} on the
 integer vectors $3\Zy_h$ and $3\Zz_{ij}$.
 \step{s4c7mem}{The vectors lie in model (i).}
 \begin{proof}\pf\ $\Zy_h\in\Phi(\typ A_8)$, and $\Zz_{ij}$ is
  $\frac13\sum_l\Ze_l-\sum_{l\in T}\Ze_l$ with $T=\{i,j,9\}$; apply
  Definition~\ref{def:3.2}, (i).~\qed\end{proof}
 \step{s4c7ip}{$(\Zy_h,\Zy_k)=1$ and $(\Zz_{ij},\Zz_{kl})=|\{i,j\}\cap\{k,l\}|$ for distinct
  vectors, and $(\Zy_h,\Zz_{ij})$ is 0 if $h\in\{i,j\}$ and 1 otherwise.}
 \begin{proof}\pf\ $(\Zy_h,\Zc)=0$ and $(\Zc,\Zc)=1$, and $(\Zc,\Ze_l)=\frac13$. Then
  $(\Zy_h,\Zy_k)=1+[h=k]$, $(\Zy_h,\Zz_{ij})=1-[h\in\{i,j\}]$ and
  $(\Zz_{ij},\Zz_{kl})=1-2+|\{i,j,9\}\cap\{k,l,9\}|$.~\qed\end{proof}
 \qedstep
 \begin{proof}\pf\ Each vector has squared norm 2 by \stepref{s4c7ip}, and the other inner
  products lie in $\{0,1\}$; so they represent a graph by Definition~\ref{def:represented}. The
  neighbours of $\Zy_1$ are the 7 vectors $\Zy_k$ and the $\binom72=21$ vectors $\Zz_{ij}$ with
  $1\notin\{i,j\}$.~\qed\end{proof}
\end{proof}

\begin{cgss}{Theorem}{4.3}\label{thm:4.3}
 \lean{SimpleGraph.IsRepresentedIn.card_le_36}\lean{SimpleGraph.IsRepresentedIn.degree_le_28}
 \lean{SimpleGraph.exists_isRepresentedIn_E₈_card_36}
 \uses{def:represented,lem:4-symsq,comp:C20,lem:4-transport,comp:C7,thm:3.5,facts:sec4}
 \assume{$X$ is represented in $\typ E_8$ by $(\Zx_u)$.}
 \prove{\begin{pfenum}
  \item $X$ has at most 36 vertices.
  \item Every vertex of $X$ has at most 28 neighbours.
  \item There is a graph represented in $\typ E_8$ with 36 vertices and a vertex with 28
   neighbours.
 \end{pfenum}}
\end{cgss}
\begin{proof}
 \pfsketch\ We will follow the paper: the Gram matrix of the symmetric squares is
 nonsingular. \repair\ The bound on the valency needs the pairing of Gap 14 (errata).
 \step{s43pd}{$4I+A(X)$ is positive definite.}
 \begin{proof}\pf\ $2I+A(X)$ is the Gram matrix of $(\Zx_u)$, so positive semidefinite by
  \path{Matrix.posSemidef_gram}; for $\Zz\ne0$,
  $\Zz^T(4I+A(X))\Zz=2|\Zz|^2+\Zz^T(2I+A(X))\Zz>0$.~\qed\end{proof}
 \step{s43gram}{The Gram matrix of $(\sigma(\Zx_u))$ in $\re^{36}$ is $4I+A(X)$.}
 \begin{proof}\pf\ Lemma~\ref{lem:4-symsq} with $\ell=8$: the entries are $(\Zx_u,\Zx_w)^2$,
  which are 4, 1 and 0 by Definition~\ref{def:represented}.~\qed\end{proof}
 \step{s43one}{Assertion 1 holds.}
 \begin{proof}\pf\ The family $(\sigma(\Zx_u))$ is linearly independent by \stepref{s43pd},
  \stepref{s43gram} and \path{Matrix.posDef_gram_iff_linearIndependent}; apply
  \path{LinearIndependent.fintype_card_le_finrank} and
  \path{finrank_euclideanSpace_fin}.~\qed\end{proof}
 \step{s43two}{Assertion 2 holds.}
 \begin{proof}\pf\ Computation~\ref{comp:C20}, for $\typ E_8$.~\qed\end{proof}
 \step{s43three}{Assertion 3 holds.}
 \begin{proof}\pf\ The graph $X_{36}$ of Computation~\ref{comp:C7} is represented in model (i);
  models (i) and (ii) are isomorphic by Theorem~\ref{thm:3.5}, so it is represented in model
  (ii) by Lemma~\ref{lem:4-transport}.~\qed\end{proof}
 \qedstep
 \begin{proof}\pf\ \stepref{s43one}, \stepref{s43two} and \stepref{s43three}.~\qed\end{proof}
\end{proof}


\begin{namedresult}{Lemma}{(trace and rank)}\label{lem:4-tracerank}
 \lean{Matrix.IsHermitian.trace_sq_le_rank_mul_sum_entry_sq}
 \uses{facts:sec4}
 \assume{$N$ is a real symmetric matrix indexed by a finite set $V$.}
 \prove{$(\tr N)^2\le\rk N\cdot\sum_{u,w\in V}N(u,w)^2$.}
 \new\ This is the inequality behind [LS73] Theorem 3.6; we will name it, since Lemma EQ1 and the
 remark after Theorem 4.4 both use it.
 \begingroup\raggedright\new\ In Lean the entrywise form stated here is
 \path{Matrix.IsHermitian.trace_sq_le_rank_mul_sum_entry_sq}, in
 \path{LineGraphs/Representation/Theorem44.lean}, which Theorem 4.4 and the remark after it use.
 The eigenvalue form $(\tr N)^2\le\rk N\cdot\sum_{u\in V}\lambda_u^2$ of the steps below is
 \path{Matrix.IsHermitian.trace_sq_le_rank_mul_sum_sq}, in
 \path{LineGraphs/Spectral/Applications.lean}; the name of the entrywise form differs, since
 that one was taken first.\par\endgroup
\end{namedresult}
\begin{proof}
 \pfsketch\ We will apply the Cauchy--Schwarz inequality to the nonzero eigenvalues.
 \pflet{$\lambda_u$, for $u\in V$, be the eigenvalues \path{Matrix.IsHermitian.eigenvalues} of
 $N$, and $Z=\{u\in V:\lambda_u\ne0\}$.}
 \step{s4trtr}{$\tr N=\sum_{u\in Z}\lambda_u$ and $\rk N=|Z|$.}
 \begin{proof}\pf\ \path{Matrix.IsHermitian.trace_eq_sum_eigenvalues}, with the zero terms
  dropped, and \path{Matrix.IsHermitian.rank_eq_card_non_zero_eigs}.~\qed\end{proof}
 \step{s4trsq}{$\sum_{u,w}N(u,w)^2=\sum_{u\in Z}\lambda_u^2$.}
 \begin{proof}\pf\ By \path{Matrix.IsHermitian.spectral_theorem}, $N=UDU^T$ with $U$ orthogonal
  and $D$ diagonal with the entries $\lambda_u$. As $N$ is symmetric,
  $\sum_{u,w}N(u,w)^2=\tr(NN^T)=\tr(UD^2U^T)=\tr(D^2U^TU)=\sum_u\lambda_u^2$, where the third
  equality is \path{Matrix.trace_mul_comm}; drop the zero terms.~\qed\end{proof}
 \qedstep
 \begin{proof}\pf\ By \path{Finset.sum_mul_sq_le_sq_mul_sq} for the functions $\lambda$ and 1 on
  $Z$, $\bigl(\sum_{u\in Z}\lambda_u\bigr)^2\le|Z|\sum_{u\in Z}\lambda_u^2$; apply
  \stepref{s4trtr} and \stepref{s4trsq}.~\qed\end{proof}
\end{proof}

\begin{namedresult}{Lemma}{(the relative bound, EQ1)}\label{lem:4-EQ1}
 \lean{InnerProductSpace.card_le_relativeBound}
 \uses{lem:4-gramfactor,lem:4-tracerank}
 \assume{The unit vectors $\Zv_1,\dots,\Zv_n$ of a real inner product space span a subspace of
 dimension $k$, where $1\le k\le7$, and $(\Zv_i,\Zv_j)=\pm\frac13$ for $i\ne j$.}
 \prove{$n\le k\cdot\frac89/(1-\frac k9)$, and in particular $n\le28$.}
\end{namedresult}
\begin{proof}
 \pfsketch\ This is [LS73] Theorem 3.6, which [CGSS] cites as [15]; we will prove it from the
 trace of the Gram matrix.
 \pflet{$M$ be the Gram matrix of the $\Zv_i$.}
 \step{s4eq1tr}{$\tr M=n$, $\tr M^2=n+\frac19n(n-1)$ and $\rk M=k$.}
 \begin{proof}\pf\ The diagonal entries of $M$ are 1 and the others $\pm\frac13$, and
  $\tr M^2=\sum_{i,j}(\Zv_i,\Zv_j)^2$; the rank is Lemma~\ref{lem:4-gramfactor},
  assertion 1.~\qed\end{proof}
 \step{s4eq1cs}{$n^2\le k\bigl(n+\frac19n(n-1)\bigr)$.}
 \begin{proof}\pf\ Lemma~\ref{lem:4-tracerank} for $M$, where
  $\sum_{i,j}M(i,j)^2=\tr M^2$ as $M$ is symmetric; apply \stepref{s4eq1tr}.~\qed\end{proof}
 \qedstep
 \begin{proof}\pf\ Divide \stepref{s4eq1cs} by $n$: $n(1-\frac k9)\le\frac89k$, with
  $1-\frac k9>0$. The bound $\frac89k/(1-\frac k9)=8k/(9-k)$ increases with $k$, and is 28 for
  $k=7$.~\qed\end{proof}
\end{proof}

\begin{namedresult}{Lemma}{(lines spanning eight dimensions, EQ2)}\label{lem:4-EQ2}
 \lean{InnerProductSpace.card_le_fourteen_of_finrank_eq_eight}
 \uses{def:represented,lem:4-rootsystem,lem:sec3-spans,comp:C20}
 \assume{$U$ is a real inner product space of dimension 8, spanned by unit vectors
 $\Zv_1,\dots,\Zv_n$ with $(\Zv_i,\Zv_j)=\pm\frac13$ for $i\ne j$.}
 \prove{$n\le14$.}
 \repair\ This is the cone construction of \path{cgss/prerequisites.md} (EQ2), which fills
 the unstated part of [LS73] (errata, Gap 7).
\end{namedresult}
\begin{proof}
 \pfsketch\ We will put the lines on a cone around a root $\Zr$ in $U\oplus\re\Ze$ and count
 the pairs of $N_\Zr$.
 \pflet{$\Zr=\sqrt2\,\Ze$ and $\Zu_i=\sqrt{3/2}\,\Zv_i+\Ze/\sqrt2$ in $U\oplus\re\Ze$, where $\Ze$
 is a unit vector orthogonal to $U$.}
 \step{s4eq2ip}{$(\Zr,\Zr)=(\Zu_i,\Zu_i)=2$, $(\Zu_i,\Zr)=1$, and
  $(\Zu_i,\Zu_j)=\frac32(\Zv_i,\Zv_j)+\frac12\in\{0,1\}$ for $i\ne j$.}
 \begin{proof}\pf\ Expand, with $(\Zv_i,\Ze)=0$.~\qed\end{proof}
 \pflet{$X'$ be the graph on $\{0,1,\dots,n\}$ in which 0 is adjacent to every $i$, and $i\ne j$
 are adjacent when $(\Zu_i,\Zu_j)=1$.}
 \step{s4eq2rep}{$\Zr,\Zu_1,\dots,\Zu_n$ represent the connected graph $X'$, and span a space of
  dimension 9.}
 \begin{proof}\pf\ Definition~\ref{def:represented} and \stepref{s4eq2ip}; 0 is adjacent to
  every other vertex. The span contains $\Ze=\Zr/\sqrt2$ and each
  $\Zv_i=(\Zu_i-\Ze/\sqrt2)/\sqrt{3/2}$, so it is $U\oplus\re\Ze$.~\qed\end{proof}
 \step{s4eq2type}{$X'$ is represented in $\typ A_9$ or in $\typ D_9$.}
 \begin{proof}\pf\ Lemma~\ref{lem:4-rootsystem}, assertion 1, with $k=9$; the models of
  $\typ E_6$, $\typ E_7$, $\typ E_8$ have ranks 6, 7, 8 by
  Lemma~\ref{lem:sec3-spans}.~\qed\end{proof}
 \qedstep
 \begin{proof}\pf\ The vertex 0 has $n$ neighbours, at most $9-1=8$ in $\typ A_9$ and at most
  $2(9-2)=14$ in $\typ D_9$, by Computation~\ref{comp:C20} and \stepref{s4eq2type}.~\qed\end{proof}
\end{proof}

\begin{namedresult}{Computation}{C8}\label{comp:C8}
 \lean{Computations.C8.quadratic_bound}
 \uses{}
 \assume{$n\le28$ and $d\le n-1$ are nonnegative integers with $n<2d+4$ and
 $0\le(d-2n+16)+14(2n-2d-4)/(2d+4-n)$.}
 \prove{$d\le16$, or $n=28$ and $d=17$.}
\end{namedresult}
\begin{proof}
 \pfsketch\ We will clear the positive denominator; the Lean proof is \texttt{decide} over the
 finitely many pairs.
 \qedstep
 \begin{proof}\pf\ Multiply by $2d+4-n>0$:
  $(d-2n+16)(2d+4-n)+14(2n-2d-4)\ge0$. Among the pairs $(n,d)$ with $n\le28$, $d\le n-1$ and
  $n<2d+4$, this holds only for $d\le16$ and for $(28,17)$; this is a finite check of the 406
  pairs with $d<n\le28$.~\qed\end{proof}
\end{proof}

\begin{namedresult}{Lemma}{(the matrices of a regular graph)}\label{lem:4-regmat}
 \lean{SimpleGraph.IsRegularOfDegree.posSemidef_centered}
 \lean{SimpleGraph.IsRegularOfDegree.rank_centered}
 \lean{SimpleGraph.IsRegularOfDegree.posSemidef_equiangular}
 \lean{SimpleGraph.IsRegularOfDegree.rank_equiangular}
 \uses{facts:sec4}
 \assume{\begin{pfenum}
  \item $X$ is $d$-regular on $n\ge1$ vertices, $Q=2I+A(X)$ is positive semidefinite, and $\Zj$
   is the all-ones vector.
  \item $t=(d+2)/(2n)$, $M=\frac12Q-tJ$ and $E=\frac23Q-\frac13J$.
 \end{pfenum}}
 \prove{\begin{pfenum}
  \item $Q\Zj=(d+2)\Zj$, and $J\Zz=0$ and $(Q\Zz,\Zj)=0$ for every $\Zz\perp\Zj$.
  \item $M\Zj=0$ and $M\Zz=\frac12Q\Zz$ for $\Zz\perp\Zj$; $M$ is positive semidefinite, and
   $\rk M=\rk Q-1$.
  \item If $n<2d+4$, then $E\Zj=\frac13(2d+4-n)\Zj$ and $E\Zz=\frac23Q\Zz$ for $\Zz\perp\Zj$; $E$
   is positive semidefinite, and $\rk E=\rk Q$.
 \end{pfenum}}
 \new\ The paper reads these facts off the eigenvalues of $A$; we will compare ranges instead, and
 state the facts once, for Theorem 4.4 and for the remark after it.
\end{namedresult}
\begin{proof}
 \pfsketch\ We will split $\re^n$ as $\re\Zj\oplus\Zj^\perp$.
 \pflet{$W=Q(\Zj^\perp)$ be the image of $\Zj^\perp$ under $Q$.}
 \step{s4rmj}{$A\Zj=d\Zj$, $Q\Zj=(d+2)\Zj$ and $J\Zj=n\Zj$.}
 \begin{proof}\pf\ \path{SimpleGraph.adjMatrix_mulVec_const_apply_of_regular} with \asm[1]{0};
  then $Q\Zj=2\Zj+A\Zj$, and the rows of $J$ sum to $n$.~\qed\end{proof}
 \step{s4rmperp}{\assume{$\Zz\perp\Zj$.}\prove{$J\Zz=0$ and $(Q\Zz,\Zj)=0$.}}
 \begin{proof}\pf\ Every entry of $J\Zz$ is $(\Zz,\Zj)=0$. Further $Q$ is symmetric, so
  $(Q\Zz,\Zj)=(\Zz,Q\Zj)=(d+2)(\Zz,\Zj)=0$ by \stepref{s4rmj}.~\qed\end{proof}
 \step{s4rmsplit}{Every $\Zy\in\re^n$ is $c\Zj+\Zz$ with $c=(\Zy,\Zj)/n$ and $\Zz\perp\Zj$.}
 \begin{proof}\pf\ $(\Zy-c\Zj,\Zj)=(\Zy,\Zj)-cn=0$, as $(\Zj,\Zj)=n\ge1$ by
  \asm[1]{0}.~\qed\end{proof}
 \step{s4rmrange}{The range of $Q$ is $\re\Zj+W$.}
 \begin{proof}\pf\ By \stepref{s4rmsplit} and \stepref{s4rmj}, $Q(c\Zj+\Zz)=c(d+2)\Zj+Q\Zz$ for
  $\Zz\perp\Zj$, and $d+2\ne0$.~\qed\end{proof}
 \step{s4rmcap}{$\re\Zj\cap W=0$.}
 \begin{proof}\pf\ If $c\Zj\in W$, then $cn=(c\Zj,\Zj)=0$ by \stepref{s4rmperp}, and $n\ge1$ by
  \asm[1]{0}; so $c=0$.~\qed\end{proof}
 \step{s4rmrk}{$\rk Q=1+\dim W$.}
 \begin{proof}\pf\ By the definition \path{Matrix.rank}, $\rk Q$ is the dimension of the range
  of $Q$, which is $\re\Zj+W$ by \stepref{s4rmrange}. Apply
  \path{Submodule.finrank_sup_add_finrank_inf_eq} with \stepref{s4rmcap}, and
  \path{finrank_span_singleton}, as $\Zj\ne0$.~\qed\end{proof}
 \step{s4rmtwo}{Assertion 2 holds.}
 \begin{proof}
  \step{s4rmMact}{$M\Zj=0$, and $M\Zz=\frac12Q\Zz$ for $\Zz\perp\Zj$.}
  \begin{proof}\pf\ $M\Zj=\frac12(d+2)\Zj-tn\Zj=0$ by \stepref{s4rmj}, since
   $tn=\frac12(d+2)$ by \asm[2]{0}; and $J\Zz=0$ by \stepref{s4rmperp}.~\qed\end{proof}
  \step{s4rmMpsd}{$M$ is positive semidefinite.}
  \begin{proof}\pf\ $M$ is symmetric, as $Q$ and $J$ are. For $\Zy=c\Zj+\Zz$ as in
   \stepref{s4rmsplit}, $(\Zy,M\Zy)=(c\Zj+\Zz,\frac12Q\Zz)=\frac12(\Zz,Q\Zz)\ge0$, by
   \stepref{s4rmMact}, \stepref{s4rmperp} and \asm[1]{0}.~\qed\end{proof}
  \step{s4rmMrk}{$\rk M=\rk Q-1$.}
  \begin{proof}\pf\ By \stepref{s4rmsplit} and \stepref{s4rmMact}, $M(c\Zj+\Zz)=\frac12Q\Zz$, so
   the range of $M$ is $W$; apply \stepref{s4rmrk} and the definition
   \path{Matrix.rank}.~\qed\end{proof}
  \qedstep
  \begin{proof}\pf\ \stepref{s4rmMact}, \stepref{s4rmMpsd} and \stepref{s4rmMrk}.~\qed\end{proof}
 \end{proof}
 \step{s4rmthree}{\assume{$n<2d+4$.}\prove{The conclusion of assertion 3 holds.}}
 \begin{proof}
  \pflet{$b=\frac13(2d+4-n)$, which is positive by \asm{1}.}
  \step{s4rmEact}{$E\Zj=b\Zj$, and $E\Zz=\frac23Q\Zz$ for $\Zz\perp\Zj$.}
  \begin{proof}\pf\ $E\Zj=\frac23(d+2)\Zj-\frac13n\Zj$ by \stepref{s4rmj}, and $J\Zz=0$ by
   \stepref{s4rmperp}.~\qed\end{proof}
  \step{s4rmEpsd}{$E$ is positive semidefinite.}
  \begin{proof}\pf\ $E$ is symmetric. For $\Zy=c\Zj+\Zz$ as in \stepref{s4rmsplit},
   $(\Zy,E\Zy)=(c\Zj+\Zz,cb\Zj+\frac23Q\Zz)=c^2bn+\frac23(\Zz,Q\Zz)\ge0$, by \stepref{s4rmEact}
   and \stepref{s4rmperp}, with $(\Zz,\Zj)=0$, $b>0$ and \asm[1]{0}.~\qed\end{proof}
  \step{s4rmErk}{$\rk E=\rk Q$.}
  \begin{proof}\pf\ By \stepref{s4rmsplit} and \stepref{s4rmEact}, $E(c\Zj+\Zz)=cb\Zj+\frac23Q\Zz$
   with $b\ne0$, so the range of $E$ is $\re\Zj+W$, which is the range of $Q$ by
   \stepref{s4rmrange}.~\qed\end{proof}
  \qedstep
  \begin{proof}\pf\ \stepref{s4rmEact}, \stepref{s4rmEpsd} and \stepref{s4rmErk}.~\qed\end{proof}
 \end{proof}
 \qedstep
 \begin{proof}\pf\ Assertion 1 is \stepref{s4rmj} and \stepref{s4rmperp}; assertions 2 and 3 are
  \stepref{s4rmtwo} and \stepref{s4rmthree}.~\qed\end{proof}
\end{proof}

\begin{cgss}{Theorem}{4.4}\label{thm:4.4}
 \lean{SimpleGraph.IsRepresentedIn.card_le_28_of_regular}
 \lean{SimpleGraph.IsRepresentedIn.degree_le_16_of_regular}
 \uses{def:represented,lem:4-gramfactor,lem:4-regmat,lem:4-sub,lem:4-symsq,lem:4-EQ1,lem:4-EQ2,lem:4-tracerank,lem:4-rootsystem,lem:sec3-spans,comp:C8,comp:C20,facts:sec4}
 \assume{\begin{pfenum}
  \item $X$ is $d$-regular on $n\ge1$ vertices.
  \item $X$ is represented in $\typ E_8$ by $(\Zx_u)$.
 \end{pfenum}}
 \prove{$n\le28$ and $d\le16$.}
 \begingroup\raggedright\new\ In Lean the theorem is proved as
 \path{SimpleGraph.IsRepresentedIn.card_le_28_of_regular_of_classification} and
 \path{SimpleGraph.IsRepresentedIn.degree_le_16_of_regular_of_classification}, in
 \path{LineGraphs/Representation/Theorem44.lean}, with the classification of the maximal
 indecomposable line systems of the nine-dimensional space of Lemma~\ref{lem:4-EQ2} as an
 explicit hypothesis. The settled names above, without the hypothesis, are proved in
 \path{LineGraphs/Representation/Theorem44Unconditional.lean}, by
 discharging it with Theorem~3.6 (\path{LineSystem.IsMaximal.exists_iso_rootType}).\par\endgroup
\end{cgss}
\begin{proof}
 \pfsketch\ We will follow the paper, with these repairs. The equiangular facts come from
 Lemmas EQ1 and EQ2 in place of [15] (errata, Error 8 and Gap 7, decision 1 of
 \path{cgss/prerequisites.md}). \repair\ The case $n=28$, $d=12$ needs its own treatment (errata,
 Error 15), and the final valency bound needs the pairing of Gap 14 (errata). The printed Gram
 matrix is that of the vectors $\sigma(\Zq_u)$, not of projections (errata, Error 13, item 6), and
 ``each eigenvalue $\alpha\ne d$'' concerns the action on $\Zj^\perp$ (item 7); we will use that
 action through quadratic forms, without counting eigenvalues.
 \pflet{$Q=2I+A(X)$, $\Zj$ the all-ones vector, $t=(d+2)/(2n)$, $M=\frac12Q-tJ$ and
 $E=\frac23Q-\frac13J$, as in Lemma~\ref{lem:4-regmat}.}
 \step{s44spec}{$Q$ is positive semidefinite of rank at most 8.}
 \begin{proof}\pf\ $Q$ is the Gram matrix of $n$ vectors of
  $\re^8$ by \asm[2]{0}, so it is positive semidefinite by \path{Matrix.posSemidef_gram}, and of rank
  at most 8 by Lemma~\ref{lem:4-gramfactor}, assertion 1, since a subspace of $\re^8$ has
  dimension at most 8 (\path{Submodule.finrank_le}).~\qed\end{proof}
 \step{s44small}{\case{$n\le17$.}\prove{$n\le28$ and $d\le16$.}}
 \begin{proof}\pf\ $d\le n-1\le16$ by \asm[1]{0} and \asm{1}.~\qed\end{proof}
 \step{s44q}{$M$ is the Gram matrix of vectors $\Zq_u\in\re^7$, with
  $(\Zq_u,\Zq_u)=1-t$, $(\Zq_u,\Zq_w)=\frac12-t$ for adjacent $u,w$, and $-t$ otherwise.}
 \begin{proof}\pf\ By Lemma~\ref{lem:4-regmat}, assertion 2, with \asm[1]{0} and
  \stepref{s44spec}, $M$ is positive semidefinite of rank $\rk Q-1\le7$. Apply
  Lemma~\ref{lem:4-gramfactor}, assertion 2, with $r=7$; the entries are those of
  $\frac12Q-tJ$.~\qed\end{proof}
 \step{s44P}{$P=(1-t)^2I+(\frac12-t)^2A+t^2(J-I-A)$ is positive semidefinite, and
  $P\Zz=(1-2t)\Zz+(\frac14-t)A\Zz$ for $\Zz\perp\Zj$.}
 \begin{proof}\pf\ By Lemma~\ref{lem:4-symsq} with $\ell=7$ and \stepref{s44q}, $P$ is the Gram
  matrix of the $\sigma(\Zq_u)\in\re^{28}$; apply \path{Matrix.posSemidef_gram}. For
  $\Zz\perp\Zj$, $J\Zz=0$ by Lemma~\ref{lem:4-regmat}, assertion 1, and $(1-t)^2-t^2=1-2t$,
  $(\frac12-t)^2-t^2=\frac14-t$.~\qed\end{proof}
 \step{s44big}{\case{$n\ge18$ and $n\ge2d+4$.}\prove{$n\le28$ and $d\le12$.}}
 \begin{proof}
  \step{s44bigpd}{$P$ is positive definite.}
  \begin{proof}\pf\ For $\Zz\perp\Zj$, $(\Zz,A\Zz)=(\Zz,Q\Zz)-2|\Zz|^2\ge-2|\Zz|^2$ by
   \stepref{s44spec}, and $\frac14-t\ge0$ by \asm{1}; so by \stepref{s44P},
   $(\Zz,P\Zz)\ge(1-2t-2(\frac14-t))|\Zz|^2=\frac12|\Zz|^2$. Further $P\Zj=p\Zj$ with
   $p=(1-t)^2+(\frac12-t)^2d+t^2(n-1-d)>0$, as $1-t\ge\frac34$ since $t\le\frac14$. $P$ is
   symmetric, being a combination of $I$, $A$ and $J$; so $(P\Zz,\Zj)=(\Zz,P\Zj)=0$ for
   $\Zz\perp\Zj$, and $P$ maps $\Zj^\perp$ into itself. So
   $(\Zy,P\Zy)=pc^2n+(\Zz,P\Zz)>0$ for $\Zy=c\Zj+\Zz\ne0$.~\qed\end{proof}
  \qedstep
  \begin{proof}\pf\ The $n$ vectors $\sigma(\Zq_u)$ are linearly independent in $\re^{28}$ by
   \stepref{s44bigpd} and \path{Matrix.posDef_gram_iff_linearIndependent}, so $n\le28$ by
   \path{LinearIndependent.fintype_card_le_finrank}; and $d\le(n-4)/2\le12$ by
   \asm{1}.~\qed\end{proof}
 \end{proof}
 \step{s44rest}{\case{$n\ge18$ and $n<2d+4$.}\prove{$n\le28$ and $d\le16$.}}
 \begin{proof}
  \pflet{$B=2(2n-2d-4)/(2d+4-n)$, and $\Zb_k=(\sqrt2\,\Zq_u(k))_{u\in V(X)}\in\re^{V(X)}$ for
   $1\le k\le7$, with the $\Zq_u$ of \stepref{s44q}.}
  \step{s44bound}{$(\Zz,Q\Zz)\le(B+2)|\Zz|^2$ for every $\Zz\perp\Zj$.}
  \begin{proof}\pf\ By \stepref{s44P}, $0\le(\Zz,P\Zz)=(1-2t)|\Zz|^2+(\frac14-t)(\Zz,A\Zz)$,
   where $\frac14-t=(n-2d-4)/(4n)<0$ by \asm{1}. Divide by it: $(\Zz,A\Zz)\le
   (1-2t)/(t-\frac14)\cdot|\Zz|^2=B|\Zz|^2$, since $1-2t=(n-d-2)/n$; add
   $2|\Zz|^2$.~\qed\end{proof}
  \step{s44cols}{$2M=\sum_{k=1}^7\Zb_k\Zb_k^T$, and $\Zb_k\perp\Zj$ for every $k$.}
  \begin{proof}\pf\ The $uw$-entry of $2M$ is $2(\Zq_u,\Zq_w)=\sum_k\Zb_k(u)\Zb_k(w)$ by
   \stepref{s44q}; and $\sum_k(\Zb_k,\Zj)^2=(\Zj,2M\Zj)=0$ by Lemma~\ref{lem:4-regmat},
   assertion 2.~\qed\end{proof}
  \step{s44ck}{$|\Zb_k|^2\le B+2$ for every $k$.}
  \begin{proof}
   \step{s44ckM}{$(\Zb_k,2M\Zb_k)=\sum_l(\Zb_l,\Zb_k)^2\ge|\Zb_k|^4$.}
   \begin{proof}\pf\ \stepref{s44cols}; drop the terms $l\ne k$.~\qed\end{proof}
   \step{s44ckQ}{$(\Zb_k,2M\Zb_k)=(\Zb_k,Q\Zb_k)\le(B+2)|\Zb_k|^2$.}
   \begin{proof}\pf\ $2M\Zb_k=Q\Zb_k$ by Lemma~\ref{lem:4-regmat}, assertion 2, as $\Zb_k\perp\Zj$ by
    \stepref{s44cols}; then \stepref{s44bound}.~\qed\end{proof}
   \qedstep
   \begin{proof}\pf\ If $\Zb_k=0$ there is nothing to prove. Otherwise divide
    $|\Zb_k|^4\le(B+2)|\Zb_k|^2$, from \stepref{s44ckM} and \stepref{s44ckQ}, by
    $|\Zb_k|^2>0$.~\qed\end{proof}
  \end{proof}
  \step{s44trace}{$0\le(d-2n+16)+7B$.}
  \begin{proof}\pf\ The diagonal entries of $2M$ are $2(1-t)$ by \stepref{s44q}, so
   $\sum_k|\Zb_k|^2=\tr(2M)=2n-d-2$ by \stepref{s44cols}; this is at most $7(B+2)$ by
   \stepref{s44ck}. \new\ The Lean proof of this step and of
   \stepref{s44bound}, in \path{LineGraphs/Representation/Theorem44.lean}, is shorter: it
   bounds the nonzero eigenvalues of $M$, at most 7 of them, through the positive semidefinite
   $P$ of \stepref{s44P}, in place of the columns $\Zb_k$.~\qed\end{proof}
  \step{s44equi}{$E$ is the Gram matrix of $n$ unit vectors $\Zw_u$ in
   $\re^8$ with $(\Zw_u,\Zw_w)=\pm\frac13$ for $u\ne w$.}
  \begin{proof}\pf\ By Lemma~\ref{lem:4-regmat}, assertion 3, with \asm[1]{0},
   \stepref{s44spec} and \asm{1}, $E$ is positive semidefinite of rank $\rk Q\le8$. Apply
   Lemma~\ref{lem:4-gramfactor}, assertion 2, with $r=8$; the entries of $E$ are 1 on the
   diagonal and $\frac23-\frac13$ or $-\frac13$ elsewhere.~\qed\end{proof}
  \step{s44n}{$n\le28$.}
  \begin{proof}\pf\ The span of the $\Zw_u$ of \stepref{s44equi} has dimension $k=\rk E\le8$
   by Lemma~\ref{lem:4-gramfactor}, assertion 1. If $k=8$, then $n\le14$ by Lemma~\ref{lem:4-EQ2}; otherwise $k\le7$ and $n\le28$ by
   Lemma~\ref{lem:4-EQ1}.~\qed\end{proof}
  \step{s44C8}{$d\le16$, or $n=28$ and $d=17$.}
  \begin{proof}\pf\ Computation~\ref{comp:C8}, with \stepref{s44trace}, \stepref{s44n} and
   \asm{1}.~\qed\end{proof}
  \step{s44rankE}{\assume{$n=28$.}\prove{The $\Zw_u$ of \stepref{s44equi} span a space of
   dimension at most 7.}}
  \begin{proof}\pf\ \repair\ This is the split at $d=12$ of Error 15 (errata): here $28<2d+4$
   by \asm{1} and \asm{2}, so $d>12$, and the case $d=12$ is in \stepref{s44big}. The span lies
   in $\re^8$, so its dimension is at most 8 by \path{Submodule.finrank_le}; it is not 8, since
   then Lemma~\ref{lem:4-EQ2} would give $n\le14$, against \asm{2}.~\qed\end{proof}
  \step{s44rankQ}{\assume{$n=28$.}\prove{$\rk Q\le7$.}}
  \begin{proof}\pf\ $\rk Q=\rk E$ by Lemma~\ref{lem:4-regmat}, assertion 3, with \asm[1]{0},
   \stepref{s44spec} and \asm{1}; and $\rk E$ is the dimension of the span of the $\Zw_u$, by
   Lemma~\ref{lem:4-gramfactor}, assertion 1, and \stepref{s44equi}. Apply \stepref{s44rankE}
   with \asm{2}.~\qed\end{proof}
  \step{s44d}{\assume{$n=28$.}\prove{$d\le16$.}}
  \begin{proof}
   \new\ The Lean proof, \path{SimpleGraph.not_regular_28_17_of_rank_le_seven}, closes the case
   $(n,d)=(28,17)$ left by \stepref{s44C8} without the components: there $t=\frac{19}{56}$, and
   $M$ has the entries $\frac{37}{56}$ on the diagonal, $\frac9{56}$ at the 17 neighbours and
   $-\frac{19}{56}$ at the other 10 vertices of each row, so $\tr M=\frac{37}2$ and
   $\sum_{u,w}M(u,w)^2=\frac{227}4$. As $\rk M=\rk Q-1\le6$ by Lemma~\ref{lem:4-regmat},
   assertion 2, and \stepref{s44rankQ}, Lemma~\ref{lem:4-tracerank} would give
   $(\frac{37}2)^2=342.25\le6\cdot\frac{227}4=340.5$, which is false. This is the argument of the
   remark after the theorem, and it needs neither Lemma~\ref{lem:4-rootsystem} nor
   Computation~\ref{comp:C20} here; we keep the route of the paper below.
   \pflet{$u$ be a vertex of $X$, $X_u$ the connected component of $X$ that contains $u$, and
    $k$ the dimension of the span of the $\Zx_w$ with $w\in V(X_u)$.}
   \step{s44dsub}{$X_u$ is connected, the $\Zx_w$ with $w\in V(X_u)$ represent it in $\re^8$, and
    $u$ has $d$ neighbours in $X_u$.}
   \begin{proof}\pf\ A component is connected and is the subgraph induced on its vertices, so
    Lemma~\ref{lem:4-sub} applies to \asm[2]{0}. The component contains every neighbour of $u$,
    and $u$ has $d$ neighbours by \asm[1]{0}.~\qed\end{proof}
   \step{s44ddim}{$k\le7$.}
   \begin{proof}\pf\ The span of the $\Zx_w$ with $w\in V(X_u)$ lies in the span of all the
    $\Zx_w$ by \path{Submodule.span_mono}, and the latter has dimension $\rk Q$ by
    Lemma~\ref{lem:4-gramfactor}, assertion 1; so $k\le\rk Q$ by \path{Submodule.finrank_mono},
    and $\rk Q\le7$ by \stepref{s44rankQ} with \asm{2}.~\qed\end{proof}
   \step{s44dtype}{$X_u$ is represented in the model of $\typ A_k$, of $\typ D_k$ with $k\ge4$,
    of $\typ E_6$ or of $\typ E_7$.}
   \begin{proof}\pf\ Lemma~\ref{lem:4-rootsystem}, assertion 1, applies by \stepref{s44dsub};
    the model of $\typ E_8$ is excluded, since it has rank $8>k$ by Lemma~\ref{lem:sec3-spans}
    and \stepref{s44ddim}.~\qed\end{proof}
   \qedstep
   \begin{proof}\pf\ By Computation~\ref{comp:C20} and \stepref{s44dtype}, $u$ has at most
    $k-1$, $2(k-2)$, 10 or 16 neighbours in $X_u$; these are at most 16 by \stepref{s44ddim}.
    Apply \stepref{s44dsub}.~\qed\end{proof}
  \end{proof}
  \qedstep
  \begin{proof}\pf\ $n\le28$ by \stepref{s44n}. If $d>16$, then $n=28$ by \stepref{s44C8}, and
   \stepref{s44d} gives $d\le16$; so $d\le16$.~\qed\end{proof}
 \end{proof}
 \qedstep
 \begin{proof}\pf\ The cases \stepref{s44small}, \stepref{s44big} and \stepref{s44rest} are
  exhaustive.~\qed\end{proof}
\end{proof}

\begin{namedresult}{Computation}{C24}\label{comp:C24}
 \lean{Computations.C24.trace_sq_gt}
 \uses{}
 \assume{$X$ is a 16-regular graph on 28 vertices, and $N=2I+A(X)-\frac9{14}J$.}
 \prove{$\tr N=38$, $\sum_{u,w}N(u,w)^2=236$, and $38^2>6\cdot236$.}
 \new\ Section 12 of \path{cgss/prerequisites.md} states C24 as $(2m-16)^2>(27-m)(192-4m)$
 for $21\le m\le27$, where $m$ is the
 multiplicity of the eigenvalue $-2$. We will count entries in its place, which needs no
 multiplicities. With Lemma (trace and rank), C24 gives only
 $\rk N\ge7$. The remark after Theorem 4.4 gets $\rk(2I+A)=8$ from this, from assertion 2 of
 Lemma (the matrices of a regular graph), which gives $\rk(2I+A)=\rk N+1$, and from the bound
 $\rk(2I+A)\le8$ of a representation in $\typ E_8$.
\end{namedresult}
\begin{proof}
 \pfsketch\ We will count the entries of each kind; the Lean proof counts them with the degrees
 and closes with \texttt{norm\_num}.
 \step{s4c24row}{Each row of $N$ has the entry $\frac{19}{14}$ on the diagonal, 16 entries
  $\frac5{14}$ and 11 entries $-\frac9{14}$.}
 \begin{proof}\pf\ The entries of $2I+A(X)$ in a row are 2 on the diagonal, 1 at the 16
  neighbours by \asm{0}, and 0 at the other 11 vertices; subtract $\frac9{14}$.~\qed\end{proof}
 \qedstep
 \begin{proof}\pf\ By \stepref{s4c24row}, $\tr N=28\cdot\frac{19}{14}=38$ and
  $\sum_{u,w}N(u,w)^2=\frac{28}{196}(361+16\cdot25+11\cdot81)=\frac{1652}7=236$; and
  $38^2=1444>1416=6\cdot236$.~\qed\end{proof}
\end{proof}

\begin{namedresult}{Remark}{(after Theorem 4.4)}\label{rem:4-after44}
 \lean{SimpleGraph.not_regular_28_16_isRepresentedIn_E₈}
 \lean{SimpleGraph.regular_28_12_and_27_16_isRepresentedIn_E₈}
 \uses{def:represented,lem:4-gramfactor,lem:4-regmat,lem:4-tracerank,lem:4-EQ2,comp:C24,lem:4-linegraph,def:glg,comp:C10,def:3.2,thm:4.4,ex:4.5,ex:4.6,facts:sec4}
 \assume{$X$ is a graph.}
 \prove{\begin{pfenum}
  \item If $X$ is 16-regular on 28 vertices, then $X$ is not represented in $\typ E_8$.
  \item The bound $n\le28$ of Theorem~\ref{thm:4.4} is attained by $L(K_8)$ and by the Chang
   graphs I, II and III, which are 12-regular on 28 vertices, and the bound $d\le16$ by the
   Schl\"afli graph, which is 16-regular on 27 vertices; all five are represented in
   $\typ E_8$.
 \end{pfenum}}
\end{namedresult}
\begin{proof}
 \pfsketch\ \repair\ The paper gives no proof of assertion 1; we will follow EQ3 of
 \path{cgss/prerequisites.md}, the repair of Gap 7 (errata). The rank of $2I+A$ is 8, so the 28
 equiangular lines of the proof of Theorem 4.4 span 8 dimensions, which Lemma EQ2 forbids.
 \new\ We will bound the rank by Lemma~\ref{lem:4-tracerank} and Computation~\ref{comp:C24}, in
 place of the count of eigenvalues in EQ3. Assertion 2 cites the examples below.
 \new\ The Lean proof, \leancite{SimpleGraph.not_regular_28_16_isRepresentedIn_E₈_of_classification}{\texttt{SimpleGraph.\allowbreak not\_\allowbreak regular\_\allowbreak 28\_\allowbreak 16\_\allowbreak isRepresentedIn\_\allowbreak E\ensuremath{{}_8}\_\allowbreak of\_\allowbreak classification}},
 runs in the other order: EQ2 first, which gives $\rk Q\le7$
 (\path{SimpleGraph.IsRepresentedIn.rank_le_seven_of_regular_of_classification}, the argument of
 the proof of Theorem~\ref{thm:4.4} for $n=28$), so that $\rk N\le6$; then C24 and
 Lemma~\ref{lem:4-tracerank} give $38^2\le6\cdot236$, which is false. It carries the same
 classification hypothesis as Theorem~\ref{thm:4.4}; the settled name, without it, is
 proved in \path{LineGraphs/Representation/Theorem44Unconditional.lean} in the same way.
 \step{s4r44one}{\assume{\begin{pfenum}
   \item $X$ is 16-regular on 28 vertices.
   \item $X$ is represented in $\typ E_8$ by $(\Zx_u)$.
  \end{pfenum}}\prove{FALSE.}}
 \begin{proof}
  \pflet{$Q=2I+A(X)$, $M=\frac12Q-\frac9{28}J$ and $E=\frac23Q-\frac13J$, the matrices of
   Lemma~\ref{lem:4-regmat} for $d=16$ and $n=28$, where $t=\frac{18}{56}=\frac9{28}$.}
  \step{s4r44psd}{$Q$ is positive semidefinite of rank at most 8.}
  \begin{proof}\pf\ $Q$ is the Gram matrix of the $\Zx_u\in\re^8$ by \asm[2]{1}, so it is
   positive semidefinite by \path{Matrix.posSemidef_gram}, and of rank at most 8 by
   Lemma~\ref{lem:4-gramfactor}, assertion 1, and \path{Submodule.finrank_le}.~\qed\end{proof}
  \step{s4r44M}{$\rk M\ge7$.}
  \begin{proof}\pf\ $2M$ is the matrix $N$ of Computation~\ref{comp:C24}, by \asm[1]{1}. By
   Lemma~\ref{lem:4-tracerank} and that computation, $38^2\le\rk(2M)\cdot236$ and
   $38^2>6\cdot236$; so $\rk M=\rk(2M)\ge7$.~\qed\end{proof}
  \step{s4r44rk}{$\rk Q=8$.}
  \begin{proof}\pf\ $\rk Q=\rk M+1\ge8$ by Lemma~\ref{lem:4-regmat}, assertion 2, with
   \asm[1]{1}, \stepref{s4r44psd} and \stepref{s4r44M}; and $\rk Q\le8$ by
   \stepref{s4r44psd}.~\qed\end{proof}
  \step{s4r44E}{$E$ is positive semidefinite of rank 8.}
  \begin{proof}\pf\ As $28<2\cdot16+4$, Lemma~\ref{lem:4-regmat}, assertion 3, applies with
   \asm[1]{1} and \stepref{s4r44psd}; it makes $E$ positive semidefinite with $\rk E=\rk Q$, and
   $\rk Q=8$ by \stepref{s4r44rk}.~\qed\end{proof}
  \step{s4r44w}{$E$ is the Gram matrix of 28 unit vectors $\Zw_u$ of $\re^8$ with
   $(\Zw_u,\Zw_w)=\pm\frac13$ for $u\ne w$.}
  \begin{proof}\pf\ Lemma~\ref{lem:4-gramfactor}, assertion 2 with $r=8$, applies by
   \stepref{s4r44E}; the entries of $E$ are 1 on the diagonal and $\frac23-\frac13$ or
   $-\frac13$ elsewhere.~\qed\end{proof}
  \step{s4r44span}{The $\Zw_u$ of \stepref{s4r44w} span a space of dimension 8.}
  \begin{proof}\pf\ Its dimension is $\rk E$ by Lemma~\ref{lem:4-gramfactor}, assertion 1, and
   \stepref{s4r44w}; and $\rk E=8$ by \stepref{s4r44E}.~\qed\end{proof}
  \qedstep
  \begin{proof}\pf\ Lemma~\ref{lem:4-EQ2}, for the span of \stepref{s4r44span} and the $\Zw_u$
   of \stepref{s4r44w}, gives $28\le14$.~\qed\end{proof}
 \end{proof}
 \step{s4r44two}{Assertion 2 holds.}
 \begin{proof}
  \step{s4r44K8}{$L(K_8)$ is 12-regular on 28 vertices and is represented in $\typ E_8$.}
  \begin{proof}\pf\ $K_8$ has $\binom82=28$ edges, and each has degree $7+7-2=12$ in
   $L(K_8)\cong L(K_8;0,\dots,0)$, by Computation~\ref{comp:C10} and Definition~\ref{def:glg}.
   Lemma~\ref{lem:4-linegraph} gives the representation by the $\Ze_i+\Ze_j$, which lie in
   $\Phi(\typ D_8)\sbs\Phi(\typ E_8)$ by Definition~\ref{def:3.2}, (ii).~\qed\end{proof}
  \qedstep
  \begin{proof}\pf\ Example~\ref{ex:4.5} gives the Chang graphs, with parameters
   $(28,12,6,4)$, and Example~\ref{ex:4.6} the Schl\"afli graph, with parameters $(27,16,10,8)$,
   each represented in $\typ E_8$; add \stepref{s4r44K8}.~\qed\end{proof}
 \end{proof}
 \qedstep
 \begin{proof}\pf\ \stepref{s4r44one} and \stepref{s4r44two}.~\qed\end{proof}
\end{proof}

The paper introduces the graphs of Examples 4.5 to 4.9 as strongly regular graphs ``which are
not line graphs''. We will prove this in the proof of Remark 4.14, as Corollary (the seven
graphs are not line graphs), since the argument uses the properties of the seven graphs that we
establish there.
We will use model (ii) of Definition 3.2 and the vector
$\Zc=\frac12\sum_{k=1}^8\Ze_k$ of $\Phi(\typ E_8)$. For a graph $X$ and $Y\sbs V(X)$, the
\emph{switched graph} $X^Y$ keeps the pairs inside $Y$ and inside $V(X)\setminus Y$, and
complements the pairs between them (SW1 of \path{cgss/prerequisites.md}).

\begin{namedresult}{Lemma}{(switching by vectors, SW4)}\label{lem:4-switch}
 \lean{SimpleGraph.IsRepresentation.switch}
 \uses{def:2.1,def:represented,def:line-system}
 \assume{\begin{pfenum}
  \item $X$ is represented in a star-closed line system $S$ by $(\Zx_u)$.
  \item $\Zh\in\Phi(S)$, $(\Zh,\Zx_u)=1$ for every $u$, and $Y\sbs V(X)$.
 \end{pfenum}}
 \prove{The family with $\Zh-\Zx_u$ for $u\in Y$ and $\Zx_u$ for $u\notin Y$ represents $X^Y$
 in $S$.}
\end{namedresult}
\begin{proof}
 \pfsketch\ We will compute the new inner products.
 \step{s4swmem}{$\Zh-\Zx_u\in\Phi(S)$.}
 \begin{proof}\pf\ $(\Zx_u,-\Zh)=-1$ by \asm[2]{0}, so $\Zx_u-\Zh$ and its negative lie in
  $\Phi(S)$, which is star-closed by \asm[1]{0}.~\qed\end{proof}
 \step{s4swip}{For $u,w\in Y$, $(\Zh-\Zx_u,\Zh-\Zx_w)=(\Zx_u,\Zx_w)$; and for $u\in Y$, $w\notin
  Y$, $(\Zh-\Zx_u,\Zx_w)=1-(\Zx_u,\Zx_w)$.}
 \begin{proof}\pf\ Expand, with $(\Zh,\Zh)=2$ and $(\Zh,\Zx_u)=(\Zh,\Zx_w)=1$ by
  \asm[2]{0}.~\qed\end{proof}
 \qedstep
 \begin{proof}\pf\ By \stepref{s4swip}, the Gram matrix keeps the entries inside $Y$ and outside
  $Y$, and replaces $0,1$ by $1,0$ between them; that is $2I+A(X^Y)$. The vectors lie in $\Phi(S)$
  by \stepref{s4swmem}.~\qed\end{proof}
\end{proof}

\begin{namedresult}{Computation}{C9}\label{comp:C9}
 \lean{SimpleGraph.changI_isSRGWith}\lean{SimpleGraph.changII_isSRGWith}
 \lean{SimpleGraph.changIII_isSRGWith}\lean{SimpleGraph.schlafli_isSRGWith}
 \lean{SimpleGraph.shrikhande_isSRGWith}\lean{SimpleGraph.clebsch_isSRGWith}
 \lean{SimpleGraph.petersen_isSRGWith}
 \lean{SimpleGraph.clebsch.foldedFiveCubeComplIso}\lean{SimpleGraph.schlafli.gq24ComplIso}
 \lean{SimpleGraph.petersen.kneserIso}
 \lean{SimpleGraph.clebsch_isSwitchingEquivalent_shrikhande}
 \lean{SimpleGraph.clebsch_isSwitchingEquivalent_lineGraph_completeBipartite}
 \uses{def:represented,def:414-folded,def:414-gq24}
 \assume{The graphs of Examples 4.5 to 4.9 are given by the integer vectors $2\Zx_u$ of their
 representations. The folded 5-cube $\square_5$ is the graph on the subsets of
 $\{1,\dots,5\}$ of even size, adjacent when their symmetric difference has 4 elements, and the
 double-six graph $Q_{27}$ is the graph on the symbols $a_i$, $b_i$ and $c_{ij}=c_{ji}$
 ($i\ne j$ in $\{1,\dots,6\}$) with $a_i\sim b_j$ for $i\ne j$, $a_i\sim c_{jl}$ and
 $b_i\sim c_{jl}$ for $i\in\{j,l\}$, and $c_{ij}\sim c_{lm}$ for disjoint $\{i,j\}$, $\{l,m\}$;
 these are \hyperref[def:414-folded]{the folded 5-cube} and \hyperref[def:414-gq24]{the
 double-six graph} of the proof of Remark 4.14.}
 \prove{The Chang graphs are strongly regular with parameters $(28,12,6,4)$, the Schl\"afli
 graph $(27,16,10,8)$, the Shrikhande graph $(16,6,2,2)$, the Clebsch graph $(16,10,6,6)$ and
 the Petersen graph of Example 4.9 $(10,3,0,1)$; the last is isomorphic to the Kneser graph of
 the 2-subsets of a 5-set; the Clebsch graph is switching equivalent to $L(K_{4,4})$ and to the
 Shrikhande graph. The Clebsch graph is isomorphic to $\comp{\square_5}$, and the Schl\"afli
 graph to $\comp{Q_{27}}$.}
\end{namedresult}
\begin{proof}
 \pfsketch\ These are finite checks on at most 28 vertices; the Lean proofs are by
 \texttt{decide} on the adjacency matrices $A(u,w)=\frac14(2\Zx_u,2\Zx_w)$, with explicit
 isomorphisms and switching sets as certificates. The two isomorphisms of the last sentence
 are the certificates of \path{cgss/proofs/clebsch/clebsch.tex}, Section 6, found by SageMath; in Lean
 they are definitions, checked by \texttt{decide}.
 \qedstep
 \begin{proof}\pf\ For each graph, \texttt{decide} checks the degree of every vertex and the
  number of common neighbours of every pair; for the Petersen graph it checks a given
  bijection, for the Clebsch graph given switching sets and bijections, and for the Clebsch and
  Schl\"afli graphs given bijections onto $\comp{\square_5}$ and $\comp{Q_{27}}$.~\qed\end{proof}
\end{proof}

\begin{cgss}{Example}{4.5}\label{ex:4.5}
 \lean{SimpleGraph.changI_isRepresentedIn_E₈}\lean{SimpleGraph.changII_isRepresentedIn_E₈}
 \lean{SimpleGraph.changIII_isRepresentedIn_E₈}
 \uses{lem:4-linegraph,lem:4-switch,comp:C9,def:3.2,thm:3.5}
 \assume{$T(8)=L(K_8)$ is represented by the vectors $\Ze_i+\Ze_j$, and the switching sets are
 $Y_{\rm I}=\{\Ze_i+\Ze_{4+i}:1\le i\le4\}$, $Y_{\rm II}=\{\Ze_i+\Ze_{i+1}:1\le i\le8\}$ (indices
 modulo 8) and $Y_{\rm III}=\{\Ze_1+\Ze_2,\Ze_2+\Ze_3,\Ze_3+\Ze_1\}\cup\{\Ze_4+\Ze_5,\Ze_5+\Ze_6,
 \Ze_6+\Ze_7,\Ze_7+\Ze_8,\Ze_8+\Ze_4\}$.}
 \prove{Replacing $\Ze_i+\Ze_j$ by $\Zc-\Ze_i-\Ze_j$ on $Y_{\rm I}$, $Y_{\rm II}$ or $Y_{\rm III}$
 gives vectors of $\Phi(\typ E_8)$ representing the Chang graphs I, II and III, which are
 strongly regular with parameters $(28,12,6,4)$.}
\end{cgss}
\begin{proof}
 \pfsketch\ We will switch with $\Zh=\Zc$.
 \step{s45base}{$T(8)$ is represented in $\typ E_8$ by the $\Ze_i+\Ze_j$, and
  $(\Zc,\Ze_i+\Ze_j)=1$.}
 \begin{proof}\pf\ Lemma~\ref{lem:4-linegraph}; the vectors lie in $\Phi(\typ D_8)\sbs
  \Phi(\typ E_8)$ by Definition~\ref{def:3.2}, (ii).~\qed\end{proof}
 \qedstep
 \begin{proof}\pf\ $\typ E_8$ is star-closed by Theorem~\ref{thm:3.5}; apply
  Lemma~\ref{lem:4-switch} to \stepref{s45base}, and Computation~\ref{comp:C9} for the
  parameters.~\qed\end{proof}
\end{proof}

\begin{cgss}{Example}{4.6}\label{ex:4.6}
 \lean{SimpleGraph.schlafli_isRepresentedIn_E₈}
 \uses{def:represented,def:3.2,comp:C9}
 \assume{$\Zx_{ij}=\Ze_i+\Ze_j$ for $1\le i<j\le6$, and $\Zy_i=\Zc-\Ze_i-\Ze_7$ and
 $\Zz_i=\Zc-\Ze_i-\Ze_8$ for $1\le i\le6$.}
 \prove{These 27 vectors lie in $\Phi(\typ E_8)$ and represent the Schl\"afli graph, which is
 strongly regular with parameters $(27,16,10,8)$.}
 \repair\ The paper prints $\frac13\sum\Ze_k$ in the second family, which gives squared norm
 $14/9$; we use $\Zc$ (errata, Error 2).
\end{cgss}
\begin{proof}
 \pfsketch\ We will check membership and the inner products.
 \step{s46mem}{The vectors lie in $\Phi(\typ E_8)$.}
 \begin{proof}\pf\ $\Zx_{ij}\in\Phi(\typ D_8)$; $\Zy_i$ and $\Zz_i$ have entries $\pm\frac12$ with
  exactly two negative ones, so they lie in model (ii) of Definition~\ref{def:3.2}.~\qed\end{proof}
 \step{s46cc}{For 2-subsets $A,B$ of $\{1,\dots,8\}$,
  $(\Zc-\sum_{a\in A}\Ze_a,\Zc-\sum_{b\in B}\Ze_b)=|A\cap B|$ and
  $(\Ze_i+\Ze_j,\Zc-\sum_{b\in B}\Ze_b)=1-|\{i,j\}\cap B|$.}
 \begin{proof}\pf\ Expand, with $(\Zc,\Zc)=2$ and $(\Zc,\Ze_k)=\frac12$.~\qed\end{proof}
 \step{s46ip}{Distinct vectors have inner products in $\{0,1\}$.}
 \begin{proof}\pf\ $(\Zx_{ij},\Zx_{kl})=|\{i,j\}\cap\{k,l\}|$; by \stepref{s46cc},
  $(\Zy_i,\Zy_k)=(\Zz_i,\Zz_k)=1$ for $i\ne k$, $(\Zy_i,\Zz_k)=[i=k]$, and
  $(\Zx_{ij},\Zy_k)=(\Zx_{ij},\Zz_k)=1-[k\in\{i,j\}]$.~\qed\end{proof}
 \qedstep
 \begin{proof}\pf\ By \stepref{s46mem} and \stepref{s46ip} the vectors represent a graph
  (Definition~\ref{def:represented}); Computation~\ref{comp:C9} gives its
  parameters.~\qed\end{proof}
\end{proof}

\begin{cgss}{Example}{4.7}\label{ex:4.7}
 \lean{SimpleGraph.shrikhande_isRepresentedIn_E₈}
 \lean{SimpleGraph.shrikhande.changIIndex}\lean{SimpleGraph.shrikhande.changIIndex_checked}
 \lean{SimpleGraph.shrikhande.changIEmbedding}
 \lean{SimpleGraph.shrikhande.nonempty_iso_induce_changI}
 \uses{lem:4-linegraph,lem:4-switch,lem:4-sub,comp:C9,def:3.2,thm:3.5,ex:4.5}
 \assume{$L(K_{4,4})$ is represented by the vectors $\Ze_i+\Ze_j$ with $1\le i\le4<j\le8$, and
 $Y=\{\Ze_i+\Ze_{4+i}:1\le i\le4\}$.}
 \prove{Replacing $\Ze_i+\Ze_j$ by $\Zc-\Ze_i-\Ze_j$ on $Y$ gives vectors of $\Phi(\typ E_8)$
 representing the Shrikhande graph, strongly regular with parameters $(16,6,2,2)$; it is an
 induced subgraph of the Chang graph I.}
 In Lean the last assertion is \path{SimpleGraph.shrikhande.nonempty_iso_induce_changI}: the
 Shrikhande graph is isomorphic to the subgraph of the Chang graph I induced on the sixteen
 vertices given by \path{SimpleGraph.shrikhande.changIIndex}, an induced embedding that the
 kernel checks (\texttt{decide +kernel}).
\end{cgss}
\begin{proof}
 \pfsketch\ We will switch with $\Zh=\Zc$, as in Example 4.5.
 \step{s47rep}{The switched vectors lie in $\Phi(\typ E_8)$ and represent the Shrikhande graph.}
 \begin{proof}\pf\ Lemma~\ref{lem:4-linegraph} for $K_{4,4}$ gives a representation in
  $\Phi(\typ D_8)\sbs\Phi(\typ E_8)$ (Definition~\ref{def:3.2}) with $(\Zc,\Ze_i+\Ze_j)=1$; $\typ
  E_8$ is star-closed by Theorem~\ref{thm:3.5}; apply Lemma~\ref{lem:4-switch}.~\qed\end{proof}
 \qedstep
 \begin{proof}\pf\ Computation~\ref{comp:C9} gives the parameters. The switched vectors are among
  those of the Chang graph I in Example~\ref{ex:4.5}, since $Y=Y_{\rm I}$, and a subfamily of a
  representation represents the induced subgraph (Lemma~\ref{lem:4-sub}).~\qed\end{proof}
\end{proof}

\begin{cgss}{Example}{4.8}\label{ex:4.8}
 \lean{SimpleGraph.clebsch_isRepresentedIn_E₈}
 \lean{SimpleGraph.clebsch.schlafliIndex}\lean{SimpleGraph.clebsch.schlafliIndex_checked}
 \lean{SimpleGraph.clebsch.schlafliEmbedding}
 \lean{SimpleGraph.clebsch.nonempty_iso_induce_schlafli}
 \uses{def:represented,lem:4-sub,def:3.2,comp:C9,ex:4.6}
 \assume{The vectors are $\Ze_i+\Ze_5$, $\Ze_i+\Ze_6$, $\Zc-\Ze_i-\Ze_7$ and $\Zc-\Ze_i-\Ze_8$,
 for $1\le i\le4$.}
 \prove{These 16 vectors are the vectors of Example 4.6 with inner product 1 with $\Ze_5+\Ze_6$;
 they represent the Clebsch graph, strongly regular with parameters $(16,10,6,6)$ and switching
 equivalent to $L(K_{4,4})$ and to the Shrikhande graph.}
 In Lean the first two assertions together are
 \path{SimpleGraph.clebsch.nonempty_iso_induce_schlafli}: the Clebsch graph is isomorphic to the
 subgraph of the Schl\"afli graph induced on the neighbours of the vertex whose vector is
 $\Ze_5+\Ze_6$, which are the vertices whose vectors have inner product 1 with it. The kernel
 checks that \path{SimpleGraph.clebsch.schlafliIndex} is an induced embedding onto these
 neighbours.
\end{cgss}
\begin{proof}
 \pfsketch\ We will read off the neighbours of $\Ze_5+\Ze_6$ in Example 4.6.
 \step{s48nbr}{The vectors of Example~\ref{ex:4.6} with inner product 1 with $\Ze_5+\Ze_6$ are
  the 16 listed.}
 \begin{proof}\pf\ Expand: $(\Ze_5+\Ze_6,\Ze_k+\Ze_l)=|\{5,6\}\cap\{k,l\}|$, which is 1 exactly for $\{k,l\}=\{i,5\}$ or
  $\{i,6\}$ with $i\le4$; and $(\Ze_5+\Ze_6,\Zc-\Ze_k-\Ze_7)=1-|\{5,6\}\cap\{k\}|$, which is 1
  exactly for $k\le4$, and the same for $\Ze_8$.~\qed\end{proof}
 \qedstep
 \begin{proof}\pf\ A subfamily of the representation of Example~\ref{ex:4.6} represents the
  induced subgraph (Lemma~\ref{lem:4-sub}); Computation~\ref{comp:C9} gives the
  parameters and the switching equivalences.~\qed\end{proof}
\end{proof}

\begin{cgss}{Example}{4.9}\label{ex:4.9}
 \lean{SimpleGraph.petersen_isRepresentedIn_E₈}
 \uses{lem:4-linegraph,lem:4-switch,comp:C9,def:3.2,thm:3.5}
 \assume{$L(K_{3,3})$ together with an isolated vertex is represented by $\Ze_4+\Ze_8$ and the
 $\Ze_i+\Ze_j$ with $1\le i\le3$ and $5\le j\le7$, and $Y=\{\Ze_i+\Ze_{4+i}:1\le i\le3\}$.}
 \prove{Replacing $\Ze_i+\Ze_j$ by $\Zc-\Ze_i-\Ze_j$ on $Y$ gives vectors of $\Phi(\typ E_8)$
 representing the Petersen graph, strongly regular with parameters $(10,3,0,1)$.}
\end{cgss}
\begin{proof}
 \pfsketch\ We will switch with $\Zh=\Zc$, as in Example 4.5.
 \step{s49base}{The ten vectors represent $L(K_{3,3})$ with an isolated vertex, in
  $\Phi(\typ E_8)$, and each has inner product 1 with $\Zc$.}
 \begin{proof}\pf\ Lemma~\ref{lem:4-linegraph} for the graph with the edges $\{i,j\}$ and
  $\{4,8\}$, whose line graph is $L(K_{3,3})$ with an isolated vertex; the vectors lie in
  $\Phi(\typ D_8)\sbs\Phi(\typ E_8)$ by Definition~\ref{def:3.2}.~\qed\end{proof}
 \qedstep
 \begin{proof}\pf\ $\typ E_8$ is star-closed by Theorem~\ref{thm:3.5}; apply
  Lemma~\ref{lem:4-switch} to \stepref{s49base}, and Computation~\ref{comp:C9} for the
  parameters and the isomorphism with the Petersen graph.~\qed\end{proof}
\end{proof}

\subsection{Theorems 4.10 and 4.11, and Remark 4.12}

\begin{cgss}{Theorem}{4.10}\label{thm:4.10}
 \lean{SimpleGraph.hoffman}
 \uses{lem:4-gram,lem:4-rootsystem,lem:4-transport,thm:4.2,def:3.1,def:glg,def:represented}
 \assume{$X$ is a connected graph, and $2I+A(X)$ is positive semidefinite.}
 \prove{$X$ is a generalized line graph, or $X$ is represented in $\typ E_8$.}
 \repair\ The paper does not assume that $X$ is connected, and the theorem is false without it:
 two disjoint Petersen graphs are neither (errata, Implicit hypothesis 1). We assume only that
 the least eigenvalue is at least $-2$; the paper assumes that it is $-2$.
\end{cgss}
\begin{proof}
 \pfsketch\ ``This follows from the Theorems 3.6 and 4.2, and the inclusion $\typ
 A_n\subseteq\typ D_{n+1}$.''
 \step{s410rep}{$X$ is represented in $\typ A_\ell$ ($\ell\ge1$), in $\typ D_\ell$
  ($\ell\ge4$), or in $\typ E_8$.}
 \begin{proof}\pf\ $X$ has a representation in $\re^{|V(X)|}$ by Lemma~\ref{lem:4-gram} and
  \asm{0}; apply Lemma~\ref{lem:4-rootsystem}, assertion 2.~\qed\end{proof}
 \step{s410A}{\case{$X$ is represented in $\typ A_\ell$.}\prove{$X$ is a generalized line
  graph.}}
 \begin{proof}\pf\ $\Phi(\typ A_\ell)\sbs\Phi(\typ D_{\ell+1})$ by Definition~\ref{def:3.1} and
  its remark on the coordinate model for small $\ell$, so
  $X$ is represented in $\typ D_{\ell+1}$ by Lemma~\ref{lem:4-transport}; apply
  Theorem~\ref{thm:4.2}, assertion 1.~\qed\end{proof}
 \step{s410D}{\case{$X$ is represented in $\typ D_\ell$.}\prove{$X$ is a generalized line
  graph.}}
 \begin{proof}\pf\ Theorem~\ref{thm:4.2}, assertion 1.~\qed\end{proof}
 \qedstep
 \begin{proof}\pf\ The cases of \stepref{s410rep} are exhaustive, and in the third $X$ is
  represented in $\typ E_8$.~\qed\end{proof}
\end{proof}

\begin{namedresult}{Computation}{C10}\label{comp:C10}
 \lean{SimpleGraph.generalizedLineGraph_degree_edge}
 \lean{SimpleGraph.generalizedLineGraph_degree_petal}
 \uses{def:glg}
 \assume{$Y$ is a graph on $\{1,\dots,m\}$, $\deg(i)$ is the degree of $i$ in $Y$, and
 $a_1,\dots,a_m\ge0$ are integers.}
 \prove{In $L(Y;a_1,\dots,a_m)$ the edge $\{i,j\}$ has degree $2(a_i+a_j)+\deg(i)+\deg(j)-2$, and the
  vertex $(i,\alpha,s)$ has degree $2(a_i-1)+\deg(i)$.}
\end{namedresult}
\begin{proof}
 \pfsketch\ We will count the neighbours of each kind; the Lean proof is parametric.
 \step{s4c10edge}{The edge $\{i,j\}$ has $\deg(i)+\deg(j)-2$ neighbours among the edges and
  $2a_i+2a_j$ among the triples.}
 \begin{proof}\pf\ The edges other than $\{i,j\}$ on $i$ or on $j$ number $(\deg(i)-1)+(\deg(j)-1)$,
  since no other edge is on both. The triples adjacent to it are the $(i,\alpha,s)$ and the
  $(j,\beta,t)$, by Definition~\ref{def:glg}.~\qed\end{proof}
 \step{s4c10petal}{$(i,\alpha,s)$ has $\deg(i)$ neighbours among the edges and $2(a_i-1)$ among
  the triples.}
 \begin{proof}\pf\ Its neighbours are the edges on $i$ and the $(i,\beta,t)$ with
  $\beta\ne\alpha$, by Definition~\ref{def:glg}.~\qed\end{proof}
 \qedstep
 \begin{proof}\pf\ Add the two counts in \stepref{s4c10edge} and in
  \stepref{s4c10petal}.~\qed\end{proof}
\end{proof}

\begin{namedresult}{Lemma}{(regular connected generalized line graphs)}\label{lem:4-regglg}
 \lean{SimpleGraph.IsGeneralizedLineGraph.lineGraph_or_cocktailParty}
 \uses{def:glg,def:CP,comp:C10}
 \assume{$X\cong L(Y;a_1,\dots,a_m)$ is connected and regular, with $Y$ on $\{1,\dots,m\}$
 and $\deg(i)$ the degree of $i$ in $Y$.}
 \prove{$X\cong L(Y')$ for some graph $Y'$, or $X\cong CP(a)$ for some $a$.}
\end{namedresult}
\begin{proof}
 \pfsketch\ This is the paper's first paragraph of the proof of Theorem 4.11. \repair\ The paper
 says that $Y$ is connected, which need not hold (errata, Gap 18), since vertices $i$ of $Y$
 with $\deg(i)=0$ and $a_i=0$ do not appear in $X$; we will split on the vertices with $a_i>0$
 instead.
 \step{s4rgzero}{\case{$a_i=0$ for every $i$.}\prove{$X\cong L(Y)$.}}
 \begin{proof}\pf\ $L(Y;0,\dots,0)\cong L(Y)$ by Definition~\ref{def:glg}.~\qed\end{proof}
 \step{s4rgiso}{\case{$a_i>0$ and $\deg(i)=0$ for some $i$.}\prove{$X\cong CP(a_i)$.}}
 \begin{proof}\pf\ No edge of $Y$ contains $i$, so by Definition~\ref{def:glg} the $2a_i\ge2$
  triples $(i,\alpha,s)$ have all their neighbours among themselves. Since $X$ is connected,
  they are all its vertices, and they span a copy of $CP(a_i)$.~\qed\end{proof}
 \step{s4rgedge}{\case{$a_i>0$ and $\deg(i)>0$ for some $i$.}\prove{This case does not occur.}}
 \begin{proof}\pf\ Pick an edge $\{i,j\}$. By regularity and C10 (Computation~\ref{comp:C10}),
  $2(a_i+a_j)+\deg(i)+\deg(j)-2=2(a_i-1)+\deg(i)$, so $2a_j+\deg(j)=0$; but $\deg(j)\ge1$.~\qed\end{proof}
 \qedstep
 \begin{proof}\pf\ Either every $a_i$ is 0, or some $a_i>0$, and then $\deg(i)=0$ or $\deg(i)>0$; so the
  three cases are exhaustive.~\qed\end{proof}
\end{proof}

\begin{cgss}{Theorem}{4.11}\label{thm:4.11}
 \lean{SimpleGraph.hoffman_rayChaudhuri}
 \uses{thm:4.10,lem:4-regglg,def:glg,def:CP,def:represented}
 \assume{$X$ is a connected regular graph, and $2I+A(X)$ is positive semidefinite.}
 \prove{$X\cong L(Y)$ for some graph $Y$, or $X\cong CP(m)$ for some $m$, or $X$ is represented
 in $\typ E_8$.}
 \new\ As in Theorem 4.10, we assume only that the least eigenvalue is at least $-2$.
\end{cgss}
\begin{proof}
 \pfsketch\ We will apply Theorem 4.10, and then the lemma on regular generalized line graphs.
 \step{s411hoff}{$X$ is a generalized line graph, or $X$ is represented in $\typ E_8$.}
 \begin{proof}\pf\ Theorem~\ref{thm:4.10}, with \asm{0}.~\qed\end{proof}
 \step{s411glg}{\case{$X\cong L(Y;a_1,\dots,a_m)$.}\prove{$X\cong L(Y')$ for some $Y'$ or
  $X\cong CP(a)$ for some $a$.}}
 \begin{proof}\pf\ Lemma~\ref{lem:4-regglg}, with \asm{0} and \asm{1}.~\qed\end{proof}
 \qedstep
 \begin{proof}\pf\ The cases of \stepref{s411hoff} are exhaustive, by
  Definition~\ref{def:glg}.~\qed\end{proof}
\end{proof}

\begin{namedresult}{Definition}{(claws)}\label{def:claw}
 \lean{SimpleGraph.IsClaw}\lean{SimpleGraph.IsFourClaw}
 \lean{SimpleGraph.ClawFree}\lean{SimpleGraph.FourClawFree}
 \uses{}
 \assume{$X$ is a graph and $k\ge1$.}
 \pflet{A \emph{$k$-claw} of $X$ with \emph{centre} $u$ is a set of $k$ pairwise non-adjacent
 neighbours of $u$; with $u$, they induce a copy of $K_{1,k}$. The graph $X$ is \emph{claw-free}
 if it has no 3-claw, and \emph{four-claw-free} if it has no 4-claw.}
 In Lean, \texttt{X.IsClaw u T} and \texttt{X.IsFourClaw u T} will say that the finite set $T$ is
 a 3-claw, or a 4-claw, of $X$ with centre $u$. Following \path{cgss/notation.md}, the
 predicates \texttt{ClawFree} and \texttt{FourClawFree} will say instead that there is no induced
 copy of the star graph \path{SimpleGraph.starGraph} on 4 or 5 vertices; the lemma below proves
 that the two readings agree. A 4-claw contains a 3-claw with the same centre.
\end{namedresult}

\begin{namedresult}{Lemma}{(claws and stars)}\label{lem:4-clawstar}
 \lean{SimpleGraph.clawFree_iff_forall_not_isClaw}
 \lean{SimpleGraph.fourClawFree_iff_forall_not_isFourClaw}
 \lean{SimpleGraph.Iso.clawFree_iff}\lean{SimpleGraph.Iso.fourClawFree_iff}
 \uses{def:claw}
 \assume{$X$ is a graph, $k\ge1$, and $K_{1,k}$ is the star graph \path{SimpleGraph.starGraph}
 on $\{0,\dots,k\}$ with centre 0.}
 \prove{\begin{pfenum}
  \item $X$ has an induced copy of $K_{1,k}$ (\path{SimpleGraph.IsIndContained}) if and only if
   $X$ has a $k$-claw.
  \item If $X\cong X'$, then $X$ has an induced copy of $K_{1,k}$ if and only if $X'$ has one.
 \end{pfenum}}
 \new\ The Lean predicates of Definition~\ref{def:claw} forbid an induced star, while our
 arguments count claws; the Lean names are the cases $k=3$ and $k=4$.
\end{namedresult}
\begin{proof}
 \pfsketch\ We will send the centre of the star to the centre of the claw.
 \step{s4csfwd}{\assume{$f$ embeds $K_{1,k}$ into $X$ as an induced subgraph.}
  \prove{$\{f(1),\dots,f(k)\}$ is a $k$-claw of $X$ with centre $f(0)$.}}
 \begin{proof}\pf\ By \path{SimpleGraph.starGraph_adj}, each $l$ with $1\le l\le k$ is adjacent
  to 0 in $K_{1,k}$, and no two of them are adjacent. As $f$ is injective and preserves
  adjacency and non-adjacency by \asm{1}, the $f(l)$ are $k$ distinct neighbours of $f(0)$, and
  no two of them are adjacent.~\qed\end{proof}
 \step{s4csbwd}{\assume{$T=\{t_1,\dots,t_k\}$ is a $k$-claw of $X$ with centre $u$.}
  \prove{The map $f$ with $f(0)=u$ and $f(l)=t_l$ for $1\le l\le k$ embeds $K_{1,k}$ into $X$ as
  an induced subgraph.}}
 \begin{proof}\pf\ $u\notin T$, since $u$ is adjacent to every member of $T$ by \asm{1} and not to
  itself; so $f$ is injective. By \asm{1} and Definition~\ref{def:claw}, $f(0)$ is adjacent to
  every $f(l)$ with $l\ge1$, and no two $f(l)$ with $l\ge1$ are adjacent; these are the
  adjacencies of $K_{1,k}$, by \path{SimpleGraph.starGraph_adj}.~\qed\end{proof}
 \step{s4csiso}{\assume{$e$ is an isomorphism from $X$ to $X'$.}\prove{Assertion 2 holds.}}
 \begin{proof}\pf\ Compose an embedding of $K_{1,k}$ with $e$ or with its inverse, by
  \path{SimpleGraph.Iso.isIndContained}, \path{SimpleGraph.Iso.isIndContained'} and
  \path{SimpleGraph.IsIndContained.trans}.~\qed\end{proof}
 \qedstep
 \begin{proof}\pf\ Assertion 1 is \stepref{s4csfwd} and \stepref{s4csbwd}; assertion 2 is
  \stepref{s4csiso}.~\qed\end{proof}
\end{proof}

\begin{namedresult}{Lemma}{(claws in a generalized line graph)}\label{lem:4-claws}
 \lean{SimpleGraph.generalizedLineGraph_not_isClaw_petal}
 \lean{SimpleGraph.generalizedLineGraph_isClaw_edge_iff}
 \lean{SimpleGraph.generalizedLineGraph_isFourClaw_edge_iff}
 \uses{def:glg,def:claw}
 \assume{$Y$ is a graph on $\{1,\dots,m\}$, $\deg(i)$ is the degree of $i$ in $Y$,
 $a_1,\dots,a_m\ge0$ are integers, and $X=L(Y;a_1,\dots,a_m)$.}
 In the Lean names, as in those of Computation C10, a \emph{petal} of $X$ is a vertex
 $(i,\alpha,s)$ of the cocktail party graph attached at $i$; it is not the petal of the proof of
 Theorem 4.2, which is a coordinate.
 \prove{\begin{pfenum}
  \item No vertex $(i,\alpha,s)$ of $X$ is the centre of a 3-claw.
  \item The edge $\{i,j\}$ of $Y$ is the centre of a 3-claw of $X$ if and only if $a_i\ge1$ and
   $a_j+\deg(j)\ge2$, or $a_j\ge1$ and $a_i+\deg(i)\ge2$.
  \item The edge $\{i,j\}$ of $Y$ is the centre of a 4-claw of $X$ if and only if $a_i\ge1$ and
   $a_j\ge1$.
 \end{pfenum}}
\end{namedresult}
\begin{proof}
 \pfsketch\ We will split the neighbours of an edge into those at its two ends.
 \pflet{For an edge $e=\{i,j\}$ of $Y$, $N_i(e)$ is the set of the edges $f\ne e$ of $Y$ with
 $i\in f$ together with the triples $(i,\alpha,s)$, and $N_j(e)$ is defined in the same way.}
 \step{s4clnbr}{The neighbours of $e=\{i,j\}$ in $X$ are the members of $N_i(e)\cup N_j(e)$,
  and $|N_j(e)|=\deg(j)-1+2a_j$.}
 \begin{proof}\pf\ By Definition~\ref{def:glg}, $e$ is adjacent to the edges that share an end
  with it and to the triples $(l,\alpha,s)$ with $l\in e$. The edges $f\ne e$ on $j$ number
  $\deg(j)-1$, and the triples $(j,\alpha,s)$ number $2a_j$.~\qed\end{proof}
 \step{s4clside}{\assume{$T\sbs N_i(e)$ is a set of pairwise non-adjacent vertices with
  $|T|\ge2$.}\prove{$T=\{(i,\alpha,0),(i,\alpha,1)\}$ for some $\alpha$; in particular
  $a_i\ge1$.}}
 \begin{proof}\pf\ By Definition~\ref{def:glg}, two edges on $i$ are adjacent, a triple
  $(i,\alpha,s)$ is adjacent to every edge on $i$, and $(i,\alpha,s)$ and $(i,\beta,r)$ are
  adjacent when $\alpha\ne\beta$. So $T$ consists of triples with one value of $\alpha$, and
  there are two of these.~\qed\end{proof}
 \step{s4clcross}{A triple $(i,\alpha,s)$ is adjacent to no member of $N_j(e)$.}
 \begin{proof}\pf\ An edge $f\ne e$ on $j$ does not contain $i$, since $e$ is the only edge on
  both $i$ and $j$; and a triple $(j,\beta,r)$ has $j\ne i$. Apply
  Definition~\ref{def:glg}.~\qed\end{proof}
 \step{s4clone}{Assertion 1 holds.}
 \begin{proof}\pf\ By Definition~\ref{def:glg}, the neighbours of $(i,\alpha,s)$ are the edges on
  $i$, each adjacent to all the others, and the triples $(i,\beta,r)$ with $\beta\ne\alpha$. So a
  3-claw would consist of three triples $(i,\beta,r)$ with one value of $\beta$, and there are
  only two of these.~\qed\end{proof}
 \step{s4cltwo}{Assertion 2 holds.}
 \begin{proof}
  \step{s4cltwoif}{\assume{$a_i\ge1$ and $a_j+\deg(j)\ge2$.}\prove{$e=\{i,j\}$ is the centre of a
   3-claw.}}
  \begin{proof}\pf\ Since $\deg(j)\ge1$, \asm{2} gives $\deg(j)-1+2a_j\ge1$, so $N_j(e)$ has a
   member $z$ by \stepref{s4clnbr}. The vertices $(i,1,0)$, $(i,1,1)$ and $z$ are neighbours of
   $e$ by \stepref{s4clnbr}; the first two are not adjacent by Definition~\ref{def:glg}, and
   neither is adjacent to $z$ by \stepref{s4clcross}.~\qed\end{proof}
  \step{s4cltwoonly}{\assume{$T$ is a 3-claw with centre $e=\{i,j\}$.}\prove{$a_i\ge1$ and
   $a_j+\deg(j)\ge2$, or $a_j\ge1$ and $a_i+\deg(i)\ge2$.}}
  \begin{proof}\pf\ $T\sbs N_i(e)\cup N_j(e)$ by \stepref{s4clnbr}, so one of the two meets $T$
   in at least two members; by the symmetry of $i$ and $j$, say $N_i(e)$. By
   \stepref{s4clside}, $a_i\ge1$ and $|T\cap N_i(e)|=2$, so $N_j(e)$ meets $T$. Thus
   $\deg(j)-1+2a_j\ge1$ by \stepref{s4clnbr}, which gives $a_j+\deg(j)\ge2$, as
   $\deg(j)\ge1$.~\qed\end{proof}
  \qedstep
  \begin{proof}\pf\ \stepref{s4cltwoif}, with $i$ and $j$ exchanged for the second alternative,
   and \stepref{s4cltwoonly}.~\qed\end{proof}
 \end{proof}
 \step{s4clthree}{Assertion 3 holds.}
 \begin{proof}
  \step{s4clthreeif}{\assume{$a_i\ge1$ and $a_j\ge1$.}\prove{$e=\{i,j\}$ is the centre of a
   4-claw.}}
  \begin{proof}\pf\ The triples $(i,1,0)$, $(i,1,1)$, $(j,1,0)$ and $(j,1,1)$ are neighbours of $e$
   by \stepref{s4clnbr}. The two at $i$ are not adjacent, nor are the two at $j$, by
   Definition~\ref{def:glg}; and those at $i$ are not adjacent to those at $j$ by
   \stepref{s4clcross}.~\qed\end{proof}
  \step{s4clthreeonly}{\assume{$T$ is a 4-claw with centre $e=\{i,j\}$.}\prove{$a_i\ge1$ and
   $a_j\ge1$.}}
  \begin{proof}\pf\ $T\sbs N_i(e)\cup N_j(e)$ by \stepref{s4clnbr}, and each of $T\cap N_i(e)$ and
   $T\cap N_j(e)$ has at most two members by \stepref{s4clside}. So each has exactly two, and
   \stepref{s4clside} gives $a_i\ge1$ and $a_j\ge1$.~\qed\end{proof}
  \qedstep
  \begin{proof}\pf\ \stepref{s4clthreeif} and \stepref{s4clthreeonly}.~\qed\end{proof}
 \end{proof}
 \qedstep
 \begin{proof}\pf\ \stepref{s4clone}, \stepref{s4cltwo} and \stepref{s4clthree}.~\qed\end{proof}
\end{proof}

\begin{cgss}{Remark}{4.12}\label{rem:4.12}
 \lean{SimpleGraph.IsGeneralizedLineGraph.clawFree_iff}
 \lean{SimpleGraph.generalizedLineGraph_fourClawFree_iff}
 \lean{SimpleGraph.lineGraph_clawFree}\lean{SimpleGraph.cocktailParty_clawFree}
 \lean{SimpleGraph.generalizedLineGraph_clawFree_iff}
 \lean{SimpleGraph.hoffman_of_clawFree}
 \uses{def:claw,lem:4-clawstar,lem:4-claws,def:glg,def:CP,thm:4.3,thm:4.10}
 \assume{$X$ is a graph.}
 \prove{\begin{pfenum}
  \item Line graphs and cocktail party graphs are claw-free.
  \item Let $Y$ be a graph on $\{1,\dots,m\}$, with $\deg(i)$ the degree of $i$ in $Y$, and let
   $a_1,\dots,a_m\ge0$ be integers. Then $L(Y;a_1,\dots,a_m)$ is claw-free if and only if $a_j=0$
   and $\deg(j)=1$ for every vertex $i$ of $Y$ with $a_i\ge1$ and every neighbour $j$ of $i$.
  \item If $X$ is a connected generalized line graph, then $X$ is claw-free if and only if
   $X\cong L(Y')$ for some graph $Y'$ or $X\cong L(K_{1,s};t,0,\dots,0)$ for some $s,t\ge0$,
   where the vertex 1 of $K_{1,s}$ is its centre. The graph $L(K_{1,s};t,0,\dots,0)$ is
   $K_{s+2t}$ with $t$ pairwise disjoint edges deleted.
  \item If $X$ is connected and claw-free, $2I+A(X)$ is positive semidefinite, and $X$ has more
   than 36 vertices or a vertex with more than 28 neighbours, then $X\cong L(Y')$ for some graph
   $Y'$ or $X\cong L(K_{1,s};t,0,\dots,0)$ for some $s,t\ge0$.
  \item For $Y$ and $a_1,\dots,a_m$ as in assertion 2, $L(Y;a_1,\dots,a_m)$ is four-claw-free if
   and only if $a_i=0$ or $a_j=0$ for every edge $\{i,j\}$ of $Y$.
 \end{pfenum}}
 \new\ The paper's $L(K_{1,n};a,0,\dots,0)$ is our $L(K_{1,s};t,0,\dots,0)$, since $n$ and $a$
 are spoken for (\path{cgss/notation.md}), and its ``sufficiently large'' is made precise by
 Theorem 4.3. The remark on Ray-Chaudhuri's theorem is history, and needs no proof.
\end{cgss}
\begin{proof}
 \pfsketch\ \repair\ The paper says ``it is easy to see''; these are the full arguments that
 Gap 9 (errata) asks for. We will read everything off the lemma on claws.
 \step{s412two}{Assertion 2 holds.}
 \begin{proof}\pf\ Claw-freeness is the absence of 3-claws, by Definition~\ref{def:claw} and, for
  the Lean predicate, Lemma~\ref{lem:4-clawstar}, assertion 1. By Lemma~\ref{lem:4-claws},
  assertion 1, only edges of $Y$ can be centres of 3-claws, and by its assertion 2 the edge $\{i,j\}$ is not one exactly when $a_i\ge1$ implies
  $a_j+\deg(j)\le1$, and the same with $i$ and $j$ exchanged. As $\deg(j)\ge1$,
  $a_j+\deg(j)\le1$ means that $a_j=0$ and $\deg(j)=1$.~\qed\end{proof}
 \step{s412one}{Assertion 1 holds.}
 \begin{proof}\pf\ Claw-freeness is preserved by isomorphisms, by
  Lemma~\ref{lem:4-clawstar}, assertion 2. A line graph $L(Y')$, with the
  vertices of $Y'$ numbered, is isomorphic to $L(Y';0,\dots,0)$ by Definition~\ref{def:glg}, and
  the condition of \stepref{s412two} holds for it trivially. By Definitions~\ref{def:CP}
  and~\ref{def:glg}, $CP(t)$ is isomorphic to $L(Y_1;t)$ for the graph $Y_1$ with one vertex and
  no edges, and the condition holds trivially again.~\qed\end{proof}
 \step{s412three}{\assume{$X$ is connected, and $X\cong L(Y;a_1,\dots,a_m)$ for $Y$ and
  $a_1,\dots,a_m$ as in assertion 2.}\prove{$X$ is claw-free if and only if $X\cong L(Y')$ for some graph $Y'$
  or $X\cong L(K_{1,s};t,0,\dots,0)$ for some $s,t\ge0$.}}
 \begin{proof}
  \step{s412ifK}{$L(K_{1,s};t,0,\dots,0)$ is claw-free.}
  \begin{proof}\pf\ The only vertex of $K_{1,s}$ with a nonzero number attached is its centre 1,
   and every neighbour of 1 has number 0 and degree 1; apply \stepref{s412two}.~\qed\end{proof}
  \step{s412only}{\assume{$X$ is claw-free, and $a_i\ge1$ for some $i$.}\prove{$X\cong
   L(K_{1,s};t,0,\dots,0)$ with $s=\deg(i)$ and $t=a_i$.}}
  \begin{proof}
   \pflet{$U$ be the set of the triples $(i,\alpha,r)$ and of the edges of $Y$ on $i$.}
   \step{s412leaf}{Every neighbour $j$ of $i$ in $Y$ has $a_j=0$ and $\deg(j)=1$.}
   \begin{proof}\pf\ $X$ is claw-free by \asm{2} and $X\cong L(Y;a_1,\dots,a_m)$ by \asm{1}, so
    $L(Y;a_1,\dots,a_m)$ is claw-free by Lemma~\ref{lem:4-clawstar}, assertion 2. Apply
    \stepref{s412two}, with \asm{2}.~\qed\end{proof}
   \step{s412closed}{Every neighbour in $L(Y;a_1,\dots,a_m)$ of a member of $U$ lies in $U$.}
   \begin{proof}\pf\ By Definition~\ref{def:glg}, the neighbours of a triple $(i,\alpha,r)$ are
    edges on $i$ and triples at $i$. The neighbours of an edge $\{i,j\}$ are the edges on $i$ or
    on $j$ and the triples at $i$ or at $j$; by \stepref{s412leaf}, the only edge on $j$ is
    $\{i,j\}$, and there are no triples at $j$.~\qed\end{proof}
   \step{s412all}{$U$ is the whole vertex set of $L(Y;a_1,\dots,a_m)$.}
   \begin{proof}\pf\ $U$ is nonempty, as $a_i\ge1$, and closed under adjacency by
    \stepref{s412closed}; and $X$ is connected by \asm{1}.~\qed\end{proof}
   \qedstep
   \begin{proof}\pf\ Number the neighbours of $i$ as $j_1,\dots,j_s$. By \stepref{s412all}, the
    map sending $\{i,j_k\}$ to the edge $\{1,k+1\}$ of $K_{1,s}$, and $(i,\alpha,r)$ to
    $(1,\alpha,r)$, is a bijection onto the vertices of $L(K_{1,s};t,0,\dots,0)$. It preserves
    adjacency by Definition~\ref{def:glg}: in both graphs the edges are pairwise adjacent and
    adjacent to every triple, and two triples are adjacent exactly when their second coordinates
    differ.~\qed\end{proof}
  \end{proof}
  \qedstep
  \begin{proof}\pf\ If $X\cong L(Y')$, then $X$ is claw-free by \stepref{s412one}; if $X\cong
   L(K_{1,s};t,0,\dots,0)$, by \stepref{s412ifK}. Conversely, let $X$ be claw-free. If every
   $a_i$ is 0, then $X\cong L(Y)$ by Definition~\ref{def:glg}; otherwise apply
   \stepref{s412only}.~\qed\end{proof}
 \end{proof}
 \step{s412clique}{$L(K_{1,s};t,0,\dots,0)$ is $K_{s+2t}$ with $t$ pairwise disjoint edges
  deleted.}
 \begin{proof}\pf\ It has $s$ edges and $2t$ triples, and by Definition~\ref{def:glg} two of its
  vertices are non-adjacent exactly when they are $(1,\alpha,0)$ and $(1,\alpha,1)$ for some
  $\alpha$; these $t$ pairs are disjoint.~\qed\end{proof}
 \step{s412four}{\assume{$X$ is connected and claw-free, $2I+A(X)$ is positive semidefinite, and
  $X$ has more than 36 vertices or a vertex with more than 28 neighbours.}\prove{$X\cong L(Y')$
  for some graph $Y'$, or $X\cong L(K_{1,s};t,0,\dots,0)$ for some $s,t\ge0$.}}
 \begin{proof}\pf\ By Theorem~\ref{thm:4.10}, whose hypotheses hold by \asm{1}, $X$ is a
  generalized line graph or is represented in $\typ E_8$; \asm{1} with Theorem~\ref{thm:4.3},
  assertions 1 and 2, excludes the second. Thus $X\cong L(Y;a_1,\dots,a_m)$ for some $Y$ and
  $a_1,\dots,a_m$ (Definition~\ref{def:glg}), and \stepref{s412three} applies, as $X$ is
  connected and claw-free by \asm{1}.~\qed\end{proof}
 \step{s412five}{Assertion 5 holds.}
 \begin{proof}\pf\ Four-claw-freeness is the absence of 4-claws, by
  Lemma~\ref{lem:4-clawstar}, assertion 1. A 4-claw contains a 3-claw with the same centre, so
  no triple is the centre of a 4-claw, by Lemma~\ref{lem:4-claws}, assertion 1; apply its
  assertion 3 to the edges.~\qed\end{proof}
 \qedstep
 \begin{proof}\pf\ Assertions 1, 2 and 5 are \stepref{s412one}, \stepref{s412two} and
  \stepref{s412five}. Assertion 3 is \stepref{s412three} and \stepref{s412clique}, since a
  generalized line graph is isomorphic to some $L(Y;a_1,\dots,a_m)$ (Definition~\ref{def:glg}).
  Assertion 4 is \stepref{s412four}.~\qed\end{proof}
\end{proof}

\subsection{Theorem 4.13 and Remark 4.14}

A graph $X$ is \emph{strongly regular} if it is regular, any two adjacent vertices have
$a$ common neighbours and any two distinct non-adjacent vertices have $c$ common
neighbours, for some $a$ and $c$; this is Mathlib's \path{SimpleGraph.IsSRGWith}, which
allows complete graphs.

\begin{namedresult}{Lemma}{(the line graph case, LG1)}\label{lem:4-LG1}
 \lean{SimpleGraph.lineGraph_isSRGWith_cases}
 \uses{def:glg,comp:C10,facts:sec4}
 \assume{$Y$ is a connected graph with at least one edge, and $L(Y)$ is strongly regular.}
 \prove{$Y$ is a cycle, or $Y\cong K_n$, $Y\cong K_{h,h}$ or $Y\cong K_{1,h}$ for some $n$ or $h$.}
\end{namedresult}
\begin{proof}
 \pfsketch\ \repair\ We will fill the two steps of the paper named in Gap 9 (errata): the values
 0 and 1 of the constant, and the bipartite case, after the proof of [BCN] Theorem 3.12.4. We
 use the corrected count of Error 7 (errata).
 \step{s4lgdeg}{$Y$ is regular, or $Y$ is bipartite and its degree is constant on each part.}
 \begin{proof}
  \step{s4lgdeg2}{\assume{$j$ and $k$ are neighbours of a vertex $i$ of $Y$.}\prove{$\deg
   j=\deg k$.}}
  \begin{proof}\pf\ The edges $\{i,j\}$ and $\{i,k\}$ have degrees $\deg i+\deg j-2$ and $\deg
   i+\deg k-2$ in $L(Y)\cong L(Y;0,\dots,0)$, by Computation~\ref{comp:C10} and
   Definition~\ref{def:glg}; these are equal, as $L(Y)$ is regular by \asm{0}.~\qed\end{proof}
  \step{s4lgwalk}{\assume{$P$ is a walk of even length in $Y$ from $u$ to $w$.}\prove{$\deg
   u=\deg w$.}}
  \begin{proof}\pf\ Induction on the length of $P$: two vertices two apart on $P$ are equal, or
   are neighbours of the vertex between them and have equal degrees by
   \stepref{s4lgdeg2}.~\qed\end{proof}
  \step{s4lgnb}{\case{$Y$ is not bipartite.}\prove{$Y$ is regular.}}
  \begin{proof}\pf\ By \path{SimpleGraph.two_colorable_iff_forall_loop_even}, with
   \path{SimpleGraph.IsBipartite}, $Y$ has a closed walk $C$ of odd length at some vertex $x$. Let
   $u,w$ be vertices; as $Y$ is connected by \asm{0}, there are walks $P$ from $u$ to $w$ and
   $R$ from $u$ to $x$. The walks $P$ and $R,C,R^{-1},P$ have lengths of opposite parity, so one
   is even, and \stepref{s4lgwalk} gives $\deg u=\deg w$.~\qed\end{proof}
  \step{s4lgbc}{\case{$Y$ is bipartite.}\prove{The degree of $Y$ is constant on each part of a
   2-colouring of $Y$.}}
  \begin{proof}\pf\ \path{SimpleGraph.IsBipartite} gives a colouring of $Y$ by
   $\mathrm{Fin}\,2$; recolour it by $\mathrm{Fin}\,2\simeq\mathrm{Bool}$ with
   \path{SimpleGraph.recolorOfEquiv} and \path{finTwoEquiv}, which keeps the colour classes, as
   \path{SimpleGraph.Coloring.even_length_iff_congr} takes a colouring by $\mathrm{Bool}$. Two
   vertices of the same colour are joined by a walk, as $Y$ is connected by \asm{0}, and it has
   even length by that lemma; apply \stepref{s4lgwalk}.~\qed\end{proof}
  \qedstep
  \begin{proof}\pf\ The two cases are exhaustive.~\qed\end{proof}
 \end{proof}
 \step{s4lgbip}{\case{$Y$ is bipartite with degrees $p\ne q$ on its parts.}\prove{$Y\cong
  K_{1,h}$.}}
 \begin{proof}\pf\ Two edges on a common vertex of degree $p$ have $p-2$ common neighbours in
  $L(Y)$, since their other ends lie in one part and are not adjacent; likewise $q-2$. Both
  kinds of pairs cannot occur, as $p\ne q$; so one degree, say $q$, is 1, and a connected graph
  with such a part is a star.~\qed\end{proof}
 \step{s4lgsmall}{\case{$Y$ is regular of valency $h\le2$.}\prove{$Y\cong K_{1,1}$ or $Y$ is a
  cycle.}}
 \begin{proof}\pf\ A connected 1-regular graph is $K_2$, and a connected 2-regular graph is a
  cycle.~\qed\end{proof}
 \step{s4lgbig}{\case{$Y$ is regular of valency $h\ge3$.}\prove{$Y\cong K_n$ or
  $Y\cong K_{h,h}$.}}
 \begin{proof}
  \step{s4lgcommon}{Two edges $\{i,j\}$, $\{i,k\}$ have $h-1$ common neighbours in $L(Y)$ if
   $j\sim k$ and $h-2$ otherwise.}
  \begin{proof}\pf\ The common neighbours are the $h-2$ other edges on $i$, and $\{j,k\}$ when
   it is an edge.~\qed\end{proof}
  \step{s4lgKn}{\case{Some $j\sim k$ with a common neighbour $i$.}\prove{$Y\cong K_n$.}}
  \begin{proof}\pf\ By \stepref{s4lgcommon} and the constant $a$, every path $j,i,k$
   closes to a triangle; so neighbourhoods are cliques, and a connected such graph is
   complete.~\qed\end{proof}
  \step{s4lgtf}{\case{$Y$ is triangle-free.}\prove{$Y\cong K_{h,h}$.}}
  \begin{proof}
   \pflet{$c$ be the constant of $L(Y)$ for non-adjacent vertices: two disjoint edges are joined by exactly $c$
    edges of $Y$.}
   \step{s4lgdisj}{Every edge $\{i,j\}$ is disjoint from some edge, and $c\le2$.}
   \begin{proof}\pf\ Pick $k\sim i$ with $k\ne j$, and $l\sim k$ with $l\ne i$, as $h\ge3$. Then
    $l\ne j$, else $i,j,k$ is a triangle; so $\{k,l\}$ is disjoint from $\{i,j\}$. Of the edges joining
    $\{i,j\}$ to a disjoint edge, at most one leaves $i$ and at most one leaves $j$, by
    triangle-freeness.~\qed\end{proof}
   \step{s4lgc0}{$c\ne0$.}
   \begin{proof}\pf\ If $c=0$, two vertices of $L(Y)$ at distance 2 would have no common
    neighbour; so the connected graph $L(Y)$ is complete, and $Y$ has no two disjoint edges,
    against \stepref{s4lgdisj}.~\qed\end{proof}
   \step{s4lgc1}{$c\ne1$.}
   \begin{proof}\pf\ Fix an edge $\{i,j\}$. An edge $\{i,k\}$, $k\ne j$, and an edge
    $\{k,l\}$, $l\ne i$, give a disjoint edge $\{k,l\}$ ($l\ne j$ by triangle-freeness); with the
    same count at $j$, there are $2(h-1)^2$ such pairs. If $c=1$, each disjoint edge arises
    once, so $|E(Y)|=1+2(h-1)+2(h-1)^2$, and $|V(Y)|=2|E(Y)|/h=4h-4+2/h$ is not an integer for
    $h\ge3$.~\qed\end{proof}
   \step{s4lgc2}{\assume{$c=2$ and $\{i,j\}$ is an edge.}\prove{$Y\cong K_{h,h}$, with parts
    $N(j)$ and $N(i)$.}}
   \begin{proof}\pf\ $N(i)$, $N(j)$ are disjoint cocliques, by triangle-freeness. For
    $x\in N(j)\setminus\{i\}$ and $y\in N(i)\setminus\{j\}$, the disjoint edges $\{x,j\}$ and
    $\{i,y\}$ are joined by $\{i,j\}$ and one more edge, which can only be $\{x,y\}$. A vertex
    $k\notin N(i)\cup N(j)$ has a neighbour $l\notin\{i,j\}$, and the two edges joining
    $\{k,l\}$ to $\{i,j\}$ put $k$ in $N(i)\cup N(j)$.~\qed\end{proof}
   \qedstep
   \begin{proof}\pf\ $c\in\{0,1,2\}$ by \stepref{s4lgdisj}; apply \stepref{s4lgc0},
    \stepref{s4lgc1} and \stepref{s4lgc2}.~\qed\end{proof}
  \end{proof}
  \qedstep
  \begin{proof}\pf\ Either some path of length 2 closes to a triangle, or $Y$ is
   triangle-free.~\qed\end{proof}
 \end{proof}
 \qedstep
 \begin{proof}\pf\ By \stepref{s4lgdeg}, $Y$ is regular or bipartite with two distinct degrees;
  the cases \stepref{s4lgbip}, \stepref{s4lgsmall} and \stepref{s4lgbig} are
  exhaustive.~\qed\end{proof}
\end{proof}

\begin{cgss}{Theorem}{4.13}\label{thm:4.13}
 \lean{SimpleGraph.seidel}
 \uses{thm:4.11,lem:4-LG1,def:CP,def:represented,facts:sec4}
 \assume{\begin{pfenum}
  \item $X$ is a connected strongly regular graph, and $2I+A(X)$ is positive semidefinite.
  \item $-2$ is an eigenvalue of $X$.
 \end{pfenum}}
 \prove{$X\cong L(K_n)$, $X\cong L(K_{n,n})$ or $X\cong CP(n)$ for some $n$, or $X$ is
 represented in $\typ E_8$.}
 \new\ The two hypotheses together are the paper's ``least eigenvalue $-2$'', read literally,
 so this is the printed theorem. We state them apart because Theorems 4.10 and 4.11 read the
 phrase as the lower bound of the first hypothesis alone, and Theorem 4.13 fails under that
 reading: for $n>36$, $K_n$ is a counterexample (errata, Implicit hypothesis 2).
\end{cgss}
\begin{proof}
 \pfsketch\ We will apply Theorem 4.11, and then the line graph case LG1.
 \step{s413hrc}{$X\cong L(Y)$, or $X\cong CP(m)$, or $X$ is represented in $\typ E_8$.}
 \begin{proof}\pf\ Theorem~\ref{thm:4.11}, with \asm[1]{0}.~\qed\end{proof}
 \step{s413L}{\case{$X\cong L(Y)$.}\prove{$X\cong L(K_n)$ or $X\cong L(K_{n,n})$.}}
 \begin{proof}
  \step{s413Y}{We may take $Y$ connected, with an edge, and $Y$ is a cycle, $K_n$, $K_{h,h}$ or
   $K_{1,h}$.}
  \begin{proof}\pf\ Remove the isolated vertices of $Y$; since $X$ is connected with a vertex,
   the rest is connected with an edge. Apply Lemma~\ref{lem:4-LG1}.~\qed\end{proof}
  \step{s413Kn}{$X$ is neither a complete graph nor $C_5$, and $X$ is not a cycle $C_p$ with
   $p\ge6$.}
  \begin{proof}\pf\ $\det(2I+A)\ne0$ for $K_n$, since $2I+A=I+J$ is positive definite, and
   $\det(2I+A(C_5))=4$; so neither has the eigenvalue $-2$ of \asm[2]{0}, which is a root of
   $\det(\mu I-A)$. In $C_p$ with $p\ge6$, vertices at distance
   2 and 3 have 1 and 0 common neighbours.~\qed\end{proof}
  \qedstep
  \begin{proof}\pf\ $L(K_{1,h})\cong K_h$ is excluded by \stepref{s413Kn}. If $Y=C_p$, then
   $X\cong C_p$, and $p=4$ by \stepref{s413Kn}, since $C_3$ is complete; and
   $C_4=L(K_{2,2})$. The rest is $L(K_n)$ or $L(K_{h,h})$, by \stepref{s413Y}.~\qed\end{proof}
 \end{proof}
 \qedstep
 \begin{proof}\pf\ The cases of \stepref{s413hrc} are exhaustive, and the other two are in the
  conclusion.~\qed\end{proof}
\end{proof}

\begin{namedresult}{Remark}{(after the examples)}\label{rem:4-large}
 \lean{SimpleGraph.hoffman_of_large}\lean{SimpleGraph.hoffman_rayChaudhuri_of_large}
 \lean{SimpleGraph.seidel_of_large}
 \uses{thm:4.3,thm:4.4,thm:4.10,thm:4.11,thm:4.13,def:glg,def:CP}
 \assume{$X$ is a connected graph, and $2I+A(X)$ is positive semidefinite.}
 \prove{\begin{pfenum}
  \item If $X$ has more than 36 vertices, or a vertex with more than 28 neighbours, then $X$ is a
   generalized line graph.
  \item If $X$ is $d$-regular on $n$ vertices, with $n>28$ or $d>16$, then $X\cong L(Y)$ for some
   graph $Y$ or $X\cong CP(m)$ for some $m$.
  \item If $X$ is strongly regular of valency $k$ on $n$ vertices, with $n>28$ or $k>16$, and $-2$
   is an eigenvalue of $X$, then $X\cong L(K_m)$, $X\cong L(K_{m,m})$ or $X\cong CP(m)$ for
   some $m$.
 \end{pfenum}}
 \new\ The paper places this remark before Theorem 4.10 and speaks of ``any subsequent
 theorem''; we will state it after Theorems 4.10, 4.11 and 4.13, with their hypotheses (errata,
 Implicit hypotheses 1 and 2).
\end{namedresult}
\begin{proof}
 \pfsketch\ We will apply the three theorems, and exclude $\typ E_8$ by Theorems 4.3 and 4.4.
 \step{s4lrone}{\assume{$X$ has more than 36 vertices, or a vertex with more than 28
  neighbours.}\prove{$X$ is a generalized line graph.}}
 \begin{proof}\pf\ By Theorem~\ref{thm:4.10} and \asm{0}, $X$ is a generalized line graph or is
  represented in $\typ E_8$; Theorem~\ref{thm:4.3}, assertions 1 and 2, with \asm{1}, excludes
  the second.~\qed\end{proof}
 \step{s4lrtwo}{\assume{$X$ is $d$-regular on $n$ vertices, and $n>28$ or $d>16$.}\prove{$X\cong
  L(Y)$ for some graph $Y$ or $X\cong CP(m)$ for some $m$.}}
 \begin{proof}\pf\ By Theorem~\ref{thm:4.11}, with \asm{0} and \asm{1}, $X\cong L(Y)$, $X\cong
  CP(m)$ or $X$ is represented in $\typ E_8$; Theorem~\ref{thm:4.4} excludes the third, by
  \asm{1}, as $n\ge1$ for a connected graph, by \asm{0}.~\qed\end{proof}
 \step{s4lrthree}{\assume{$X$ is strongly regular of valency $k$ on $n$ vertices, $n>28$ or
  $k>16$, and $-2$ is an eigenvalue of $X$.}\prove{$X\cong L(K_m)$, $X\cong L(K_{m,m})$ or
  $X\cong CP(m)$ for some $m$.}}
 \begin{proof}\pf\ By Theorem~\ref{thm:4.13}, with \asm{0} and \asm{1}, $X$ is one of these or is
  represented in $\typ E_8$; Theorem~\ref{thm:4.4} excludes the last, as $X$ is $k$-regular on
  $n$ vertices by \asm{1}, and $n\ge1$ for a connected graph, by \asm{0}.~\qed\end{proof}
 \qedstep
 \begin{proof}\pf\ \stepref{s4lrone}, \stepref{s4lrtwo} and \stepref{s4lrthree}.~\qed\end{proof}
\end{proof}

\medskip\noindent\textbf{What remains of Section 4.}

Remark 4.14 is stated below, and its proof is in Section~\ref{sec:rem414}.
The remarks after Theorems 4.1 and 4.2, the last pieces of the section, are stated and proved
after those theorems.

\begin{cgss}{Remark}{4.14}\label{rem:4.14}
 \lean{SimpleGraph.seidel_seven}
 \uses{thm:4.11,def:CP,def:represented}
 \assume{$X$ is a connected strongly regular graph with least eigenvalue $-2$, represented in
 $\typ E_8$, and isomorphic to none of $L(K_m)$, $L(K_{m,m})$ and $CP(m)$.}
 \prove{$X$ is isomorphic to one of the Chang graphs I, II and III, the Schl\"afli graph, the
 Shrikhande graph, the Clebsch graph and the Petersen graph; and these seven graphs are
 pairwise non-isomorphic and have these properties.}
 \repair\ The paper gives no proof, and ``the seven graphs'' needs the Chang graphs to be
 pairwise non-isomorphic (errata, Gap 12). We will give its proof, by the
 route of \path{cgss/prerequisites.md} (Section 7.3), in Section~\ref{sec:rem414}.
\end{cgss}

\endgroup
